% Luc Destombe --- licence CC BY-NC-SA 4.0
% https://creativecommons.org/licenses/by-nc-sa/4.0/deed.fr
% Fichier autonome : compiler avec pdflatex, deux fois (signets, nombre de pages).
\documentclass[11pt,a4paper]{article}
\makeatletter
\newif\iflycee@unecolonne
\lycee@unecolonnetrue
\newif\iflycee@palatino
\lycee@palatinotrue

\RequirePackage[utf8]{inputenc}
\RequirePackage[T1]{fontenc}
\RequirePackage[french]{babel}
\RequirePackage{etoolbox}

\newcommand{\ShorthandsOff}{\@ifundefined{shorthandoff}{}{\shorthandoff{;:!?}}}
\newcommand{\ShorthandsOn}{\@ifundefined{shorthandon}{}{\shorthandon{;:!?}}}
\ShorthandsOff
\AtBeginDocument{\ShorthandsOn}
\AtBeginEnvironment{tikzpicture}{\ShorthandsOff}

\RequirePackage{geometry}
\RequirePackage{fancyhdr}
\RequirePackage{lastpage}
\RequirePackage{multicol}
\RequirePackage{array}
\RequirePackage{enumitem}
\RequirePackage{xcolor}
\RequirePackage{needspace}

\RequirePackage{amsmath,amssymb}
\RequirePackage{mathpazo}
\DeclareMathSizes{10.95}{10.95}{8}{5.5}
\begingroup\catcode`\_=\active \catcode`\^=\active
\gdef_#1{\sb{\scriptscriptstyle #1}}
\gdef^#1{\sp{\scriptscriptstyle\lycee@moinsepais #1}}
\endgroup
\newcommand\lycee@moins{\mathbin{%
  \setbox\z@\hbox{$\m@th\scriptscriptstyle\lycee@moinsorig$}%
  \dimen@=.55\wd\z@ \advance\dimen@-.5pt
  \hbox to.75\wd\z@{\hss$\m@th\scriptscriptstyle\vcenter{\hbox{%
    \tikz\draw[line width=.5pt,line cap=round](0,0)--(\the\dimen@,0);}}$\hss}}}
\begingroup\lccode`\~=`\- \lowercase{\endgroup\let~}\lycee@moins
\newcommand\lycee@moinsepais{\mathcode`\-="8000 }
\AtBeginDocument{\mathchardef\lycee@moinsorig=\mathcode`\-\relax}
\AtBeginDocument{\catcode`\_=12 \mathcode`\_="8000
  \catcode`\^=12 \mathcode`\^="8000 }
\newcommand\lycee@exposants{%
  \fontdimen13\textfont2=.48\fontdimen6\textfont2
  \fontdimen14\textfont2=.48\fontdimen6\textfont2
  \fontdimen15\textfont2=.48\fontdimen6\textfont2
  \fontdimen13\scriptfont2=.48\fontdimen6\scriptfont2
  \fontdimen14\scriptfont2=.48\fontdimen6\scriptfont2
  \fontdimen15\scriptfont2=.48\fontdimen6\scriptfont2}
\every@math@size\expandafter{\the\every@math@size\lycee@exposants}
\newlength{\retraitcalculs}
\setlength{\retraitcalculs}{2em}

\RequirePackage{tikz}
\usetikzlibrary{arrows.meta,calc,babel,shapes.misc}
\RequirePackage[most]{tcolorbox}
\tcbuselibrary{skins,breakable}

\definecolor{coulDef}{HTML}{1F4E79}
\definecolor{coulProp}{HTML}{7030A0}
\definecolor{coulRem}{HTML}{C55A11}
\definecolor{coulEx}{HTML}{2E7D32}
\definecolor{coulExo}{HTML}{B03A2E}
\definecolor{coulPart}{HTML}{1F4E79}
\definecolor{coulMeth}{HTML}{1B6E4F}
\definecolor{coulCourbe}{HTML}{0B5ED7}
\definecolor{coulCourbeB}{HTML}{C00000}
\definecolor{coulCourbeC}{HTML}{2E7D32}
\definecolor{coulGrille}{HTML}{8C8C8C}
\definecolor{coulRep}{HTML}{1B6E4F}
\definecolor{coulBar}{HTML}{2E7D32}

\newcommand{\R}{\mathbb{R}}
\newcommand{\Rs}{\mathbb{R}^{*}}
\renewcommand{\le}{\leqslant}
\renewcommand{\ge}{\geqslant}
\renewcommand{\approx}{\simeq}
\newcommand{\vect}[1]{\overrightarrow{#1}}
\newcommand{\norme}[1]{\left\lVert #1 \right\rVert}
\newcommand{\intervalle}[2]{\left[#1\,;#2\right]}
\RequirePackage{cancel}
\renewcommand{\CancelColor}{\color{coulExo}}
\newcommand{\fois}[1]{\times{\color{coulExo}#1}}
\newcommand{\barre}[1]{\cancel{#1}}

\tikzset{grille/.style={coulGrille,line width=0.4pt}}
\newcommand{\quadrillage}[4]{\draw[grille,step=1] (#1,#2) grid (#3,#4);}
\tikzset{
  croix/.style={cross out,draw,line width=0.6pt,inner sep=0pt,minimum size=4pt},
  croixdroite/.style={croix,rotate=45,line width=0.8pt},
}
\newcommand{\pt}[3]{\path (#1) node[croix]{} node[#2,font=\small,inner sep=2.5pt]{$#3$};}

\newcommand{\lycee@repere}[4]{%
  \draw[grille,step=1] (#1,#2) grid (#3,#4);
  \draw[-{Stealth[length=2mm]},thick] (#1,0)--({#3+0.4},0) node[below left=-1pt]{$x$};
  \draw[-{Stealth[length=2mm]},thick] (0,#2)--(0,{#4+0.4}) node[below left=-1pt]{$y$};
  \node[below left,font=\scriptsize,inner sep=1.5pt] at (0,0) {$0$};}
\newcommand{\lycee@gradx}[1]{\foreach \k/\l in {#1}{%
  \draw (\k,0.07)--(\k,-0.07) node[below,font=\scriptsize,inner sep=1.5pt]{$\l$};}}
\newcommand{\lycee@grady}[1]{\foreach \k/\l in {#1}{%
  \draw (0.07,\k)--(-0.07,\k) node[left,font=\scriptsize,inner sep=1.5pt]{$\l$};}}
\newcommand{\lycee@intff}[2]{\left[#1\,;#2\right]}
\newcommand{\lycee@intfo}[2]{\left[#1\,;#2\right[}
\newcommand{\lycee@intof}[2]{\left]#1\,;#2\right]}
\newcommand{\lycee@intoo}[2]{\left]#1\,;#2\right[}
\newcommand{\lycee@ens}[1]{\left\{#1\right\}}
\AtBeginDocument{%
  \providecommand{\repere}{\lycee@repere}%
  \providecommand{\gradx}{\lycee@gradx}%
  \providecommand{\grady}{\lycee@grady}%
  \providecommand{\Cf}{\mathcal{C}\sb{\scriptscriptstyle f}}%
  \providecommand{\Cg}{\mathcal{C}\sb{\scriptscriptstyle g}}%
  \providecommand{\intff}{\lycee@intff}%
  \providecommand{\intfo}{\lycee@intfo}%
  \providecommand{\intof}{\lycee@intof}%
  \providecommand{\intoo}{\lycee@intoo}%
  \providecommand{\ens}{\lycee@ens}}

\newlist{questions}{enumerate}{3}
\setlist[questions,1]{label=\arabic*\textdegree),leftmargin=*,itemsep=2pt,topsep=3pt}
\setlist[questions,2]{label=\alph*),leftmargin=*,itemsep=1pt,topsep=2pt}
\setlist[questions,3]{label=\roman*),leftmargin=*,itemsep=1pt,topsep=1pt}
\setlist[itemize]{label=\textbullet,leftmargin=*,itemsep=2pt,topsep=3pt}

\newcommand{\besoinplace}[1]{\par\penalty\@M\begingroup
  \dimen@=#1\relax\dimen@ii=\pagegoal\advance\dimen@ii-\pagetotal
  \ifdim\dimen@>\dimen@ii\ifdim\dimen@ii>\z@\vfil\fi\break\fi\endgroup}

\newcommand{\lycee@pied}{\thepage/\pageref{LastPage}}
\newcommand{\entete}[2]{%
  \pagestyle{fancy}\fancyhf{}%
  \fancyhead[L]{\footnotesize\color{coulPart}\textbf{#1}}%
  \fancyhead[R]{\footnotesize\color{coulPart}\textbf{#2}}%
  \fancyfoot[C]{\footnotesize\lycee@pied}%
  \renewcommand{\headrulewidth}{0.4pt}%
  \renewcommand{\headrule}{\hbox to\headwidth{%
    \color{coulPart!40}\leaders\hrule height \headrulewidth\hfill}}}

\RequirePackage{ccicons}
\newcommand{\lycee@licence}{{\scriptsize\color{gray}%
  \copyright\ Luc Destombe --- \ccbyncsa\ CC BY-NC-SA 4.0}}
\newif\iflycee@licence \lycee@licencetrue
\newcommand{\sanslicence}{\lycee@licencefalse}
\AtBeginDocument{\iflycee@licence\AddToHookNext{shipout/foreground}{%
  \put(0,0){\rlap{\kern\dimexpr1in+\hoffset+\oddsidemargin+\textwidth\relax
    \llap{\raisebox{-\dimexpr1in+\voffset+\topmargin+\headheight+\headsep
      +\textheight+\footskip\relax}{\lycee@licence}}}}}\fi}

\newcommand{\formulesserrees}{%
  \setlength{\abovedisplayskip}{3pt}\setlength{\belowdisplayskip}{3pt}%
  \setlength{\abovedisplayshortskip}{2pt}\setlength{\belowdisplayshortskip}{3pt}}

\setlength{\parindent}{0pt}

\RequirePackage[implicit=false,hidelinks,pdfencoding=auto,psdextra,bookmarksopen,
  bookmarksopenlevel=1,pdfcreator={},pdfproducer={}]{hyperref}
\RequirePackage{bookmark}
\hypersetup{pdfauthor={Luc Destombe},pdfkeywords={CC BY-NC-SA 4.0}}
\newcounter{lycee@signet}
\newcommand{\signet}[3][]{\stepcounter{lycee@signet}%
  \edef\lycee@dest{\if\relax\detokenize{#1}\relax
    signet.\arabic{lycee@signet}\else#1\fi}%
  \hypertarget{\lycee@dest}{}\bookmark[level=#2,dest=\lycee@dest]{#3}}
\pdfstringdefDisableCommands{%
  \def\R{R}\def\vect#1{#1}\def\Cf{Cf}\def\Cg{Cg}\def\fois#1{×#1}%
  \def\barre#1{#1}\def\intervalle#1#2{[#1;#2]}\def\intff#1#2{[#1;#2]}%
  \def\textsuperscript#1{#1}\def\dfrac#1#2{#1/#2}\def\frac#1#2{#1/#2}%
  \def\vec#1{#1⃗}}%
\RequirePackage[official]{eurosym}

\tcbset{exercice corrige/.style={
  enhanced,breakable,before upper=\raggedright,
  colback=white,colframe=coulExo,
  boxrule=0.8pt,arc=2pt,
  left=6pt,right=6pt,top=8pt,bottom=5pt,
  before skip=9pt,after skip=4pt,
  coltitle=white,fonttitle=\bfseries\small,
  attach boxed title to top left={xshift=8pt,yshift=-\tcboxedtitleheight/2},
  boxed title style={colback=coulExo,arc=2pt,boxrule=0pt}}}
\newcounter{exo}
\newtcolorbox{exercice}[1][]{exercice corrige,
  phantom={\signet[exercice-\arabic{exo}]{2}{Exercice \arabic{exo}}},
  title={Exercice \theexo\ifx\relax#1\relax\else{} --- #1\fi}}
\newenvironment{exo}[1][\relax]{\refstepcounter{exo}\begin{exercice}[#1]}{\end{exercice}}

\geometry{top=1.8cm,bottom=1cm,left=1cm,right=1cm,headheight=14pt,footskip=0.6cm}
\renewcommand{\lycee@pied}{\thepage}

\renewcommand{\theexo}{\ifnum\value{exo}<10 0\fi\arabic{exo}}
\tcbset{exercice corrige/.style={
  enhanced,breakable,before upper=\raggedright,
  colback=white,colframe=coulExo,
  boxrule=0.9pt,arc=2pt,
  left=8pt,right=8pt,top=8pt,bottom=6pt,
  before skip=8pt,after skip=6pt,
  coltitle=white,fonttitle=\bfseries,
  attach boxed title to top left={xshift=8pt,yshift=-\tcboxedtitleheight/2},
  boxed title style={colback=coulExo,arc=2pt,boxrule=0pt}}}

\newcommand{\entetefiche}[2]{%
  \begin{center}
    \begin{tcolorbox}[enhanced,width=0.92\linewidth,colback=white,colframe=coulPart,
        boxrule=1.4pt,arc=3pt,halign=center,top=6pt,bottom=6pt]
      {\large\bfseries\color{coulPart} CORRIGÉ --- FICHE D'EXERCICES}\\[3pt]
      {\normalsize\bfseries #1}\\[2pt]
      {\normalsize #2}
    \end{tcolorbox}
  \end{center}}

\newcounter{partie}
\newcommand{\partie}[1]{%
  \stepcounter{partie}%
  \besoinplace{8\baselineskip}\vspace{4pt}%
  \begin{tcolorbox}[enhanced,colback=coulPart!8,colframe=coulPart,boxrule=0pt,
      leftrule=3pt,arc=0pt,left=6pt,right=4pt,top=3pt,bottom=3pt,
      before skip=6pt,after skip=2pt]
    \bfseries\color{coulPart}\Roman{partie} --- #1\signet{1}{\Roman{partie} --- #1}
  \end{tcolorbox}}
\newcommand{\partiesansnum}[1]{%
  \besoinplace{8\baselineskip}\vspace{4pt}%
  \begin{tcolorbox}[enhanced,colback=coulPart!8,colframe=coulPart,boxrule=0pt,
      leftrule=3pt,arc=0pt,left=6pt,right=4pt,top=3pt,bottom=3pt,
      before skip=6pt,after skip=2pt]
    \bfseries\color{coulPart}#1\signet{1}{#1}
  \end{tcolorbox}}

\setlist[questions,1]{label=\arabic*\textdegree),leftmargin=*,itemsep=3pt,topsep=3pt}
\setlist[questions,2]{label=\alph*),leftmargin=*,itemsep=2pt,topsep=2pt}
\setlist[questions,3]{label=\roman*),leftmargin=*,itemsep=1pt,topsep=1pt}
\newlist{expr}{enumerate}{1}
\setlist[expr]{label={\Alph*\ $=$},leftmargin=*,itemsep=2pt,topsep=3pt}

\newcommand{\res}[1]{{\color{coulRep}#1}}
\newcommand{\rem}[1]{\par\smallskip{\small\itshape\color{black!60}$\rightarrow$ #1}}
\newenvironment{calculs}{\par\smallskip\noindent$\displaystyle\begin{aligned}[t]}%
  {\end{aligned}$\par\smallskip}
\newcommand{\justif}[1]{\quad\text{\small\itshape\color{black!60}#1}}
\newcommand{\cas}[1]{\par\smallskip\noindent\textbf{\color{coulExo}#1}\nobreak}
\newcommand{\explique}[2]{\par\smallskip\noindent
  \begin{minipage}[t]{0.6\linewidth}\raggedright #1\end{minipage}\hfill
  \begin{minipage}[t]{0.37\linewidth}\raggedright\leavevmode
    {\small\itshape\color{black!60}#2}\end{minipage}\par}

\setlength{\parskip}{2pt}

\makeatother
\geometry{top=1.8cm,bottom=1.6cm,left=1.8cm,right=1.8cm,headheight=14pt,footskip=30pt}
\entete{Chapitre 2 --- Fonctions --- corrigé}{Première spécialité}
\usepackage{tkz-tab}

\let\casligne\cas
\renewcommand{\cas}[1]{\casligne{#1}\par\nopagebreak}
\newcommand{\gris}[1]{{\small\itshape\color{black!60}#1}}
\tikzset{lecture/.style={coulProp,dashed,line width=0.7pt},
  inter/.style={red,line width=2pt},
  interB/.style={orange,line width=2pt}}
\newcommand{\courbeExQuatre}{\draw[coulCourbe,line width=1.1pt] (-3,4)
    .. controls (-2.667,3) and (-2.333,1.833) .. (-2,1)
    .. controls (-1.667,0.167) and (-1.333,-0.5) .. (-1,-1)
    .. controls (-0.667,-1.5) and (-0.333,-2) .. (0,-2)
    .. controls (0.333,-2) and (0.667,-1.5) .. (1,-1)
    .. controls (1.333,-0.5) and (1.667,0.5) .. (2,1)
    .. controls (2.333,1.5) and (2.667,2) .. (3,2)
    .. controls (3.333,2) and (3.667,1.5) .. (4,1)
    .. controls (4.333,0.5) and (4.667,-0.333) .. (5,-1);
  \path[coulCourbe] (-3,4) node[croix]{} (5,-1) node[croix]{};
  \node[coulCourbe,font=\footnotesize] at (-2.3,4.4) {$\Cg$};}
\newcommand{\courbesExDix}{\draw[coulCourbe,line width=1.1pt] (-3,3)
    .. controls (-2.667,2.333) and (-2.333,1.667) .. (-2,1)
    .. controls (-1.667,0.333) and (-1.333,-0.5) .. (-1,-1)
    .. controls (-0.667,-1.5) and (-0.333,-2) .. (0,-2)
    .. controls (0.333,-2) and (0.667,-0.667) .. (1,0)
    .. controls (1.333,0.667) and (1.667,1.5) .. (2,2)
    .. controls (2.333,2.5) and (2.667,3) .. (3,3)
    .. controls (3.333,3) and (3.667,2.5) .. (4,2)
    .. controls (4.333,1.5) and (4.667,0.667) .. (5,0);
  \draw[coulCourbeC,line width=1.1pt] (-3,-1)
    .. controls (-2.667,-0.333) and (-2.333,0.5) .. (-2,1)
    .. controls (-1.667,1.5) and (-1.333,2) .. (-1,2)
    .. controls (-0.667,2) and (-0.333,1.333) .. (0,1)
    .. controls (0.333,0.667) and (0.667,0.333) .. (1,0)
    .. controls (1.333,-0.333) and (1.667,-1) .. (2,-1)
    .. controls (2.333,-1) and (2.667,-0.5) .. (3,0)
    .. controls (3.333,0.5) and (3.667,1.5) .. (4,2)
    .. controls (4.333,2.5) and (4.667,2.667) .. (5,3);
  \node[coulCourbe,font=\footnotesize] at (2.6,3.6) {$\Cf$};
  \node[coulCourbeC,font=\footnotesize] at (5.5,3.4) {$\Cg$};}
\newcommand{\courbeExQuatorze}{\draw[coulCourbe,line width=1.1pt] (-4,0)
    .. controls (-3.667,0.333) and (-3.333,0.667) .. (-3,1)
    .. controls (-2.667,1.333) and (-2.333,2) .. (-2,2)
    .. controls (-1.667,2) and (-1.333,1.333) .. (-1,1)
    .. controls (-0.667,0.667) and (-0.333,0.5) .. (0,0)
    .. controls (0.333,-0.5) and (0.667,-1.5) .. (1,-2)
    .. controls (1.333,-2.5) and (1.667,-3) .. (2,-3)
    .. controls (2.333,-3) and (2.667,-2.5) .. (3,-2)
    .. controls (3.333,-1.5) and (3.667,-0.5) .. (4,0)
    .. controls (4.333,0.5) and (4.667,0.5) .. (5,1)
    .. controls (5.333,1.5) and (5.667,2.333) .. (6,3);
  \path[coulCourbe] (-4,0) node[croix]{} (6,3) node[croix]{};
  \node[coulCourbe,font=\footnotesize] at (5.4,3.6) {$\Cf$};}
\newcommand{\repereExQuatorze}{\repere{-5}{-4}{7}{5}
  \gradx{-4,-3,-2,-1,1,2,3,4,5,6} \grady{-3,-2,-1,1,2,3,4}}
\setlength{\parindent}{0pt}

\begin{document}

\entetefiche{Chapitre 2 --- Fonctions}{Première spécialité mathématiques}

\partie{Images et antécédents}

\begin{exo}[Vocabulaire]
\cas{1) a)}
\explique{$f(3)=\res{-1}$}{« l'image de $3$ est $-1$ » : le nombre de départ entre
parenthèses, son image après le signe $=$}
\cas{1) b)}
\explique{$g(-2)=\res{6}$}{« $-2$ est un antécédent de $6$ » : $-2$ est le nombre de
départ, $6$ est son image}
\cas{1) c)}
\explique{$h(0)=\res{5}$}{un point est sur la courbe lorsque son ordonnée est l'image de
son abscisse : ici $5$ est l'image de $0$}
\cas{1) d)}
\explique{$k(4)=\res{9}$}{« $4$ a pour image $9$ » : même phrase qu'en a), dans l'autre
ordre}
\cas{2)}
\explique{L'\res{image} de $5$ par $f$ est $-7$.}{$f(5)=-7$ : $5$ est le nombre de
départ, $-7$ son image}
\explique{Le nombre $5$ est \res{un antécédent} de $-7$ par $f$.}{on dit « un »
antécédent : $-7$ peut en avoir d'autres}
\explique{Le point de coordonnées $(5\,;-7)$ est sur la \res{courbe} de $f$.}{abscisse :
le nombre de départ ; ordonnée : son image}
\rem{On dit « \emph{l'}image » (un nombre n'a qu'une image) mais « \emph{un}
antécédent » (un nombre peut avoir plusieurs antécédents, ou aucun).}
\end{exo}

\begin{exo}[Calculer des images]
\explique{$f(x)=2x^{2}-5x$ : pour calculer une image, on remplace $x$ par le nombre,
\textbf{entre parenthèses}.}{méthode 1 ; priorités : le carré d'abord, puis les produits,
puis la soustraction}
\cas{1)}
\begin{calculs}
f(3) &= 2\times 3^{2}-5\times 3 &&\justif{on remplace $x$ par $3$}\\
f(3) &= 2\times 9-15 &&\justif{le carré avant le produit : $3^{2}=9$}\\
f(3) &= \res{3} &&
\end{calculs}
\begin{calculs}
f(-2) &= 2\times(-2)^{2}-5\times(-2) &&\justif{$-2$ entre parenthèses}\\
f(-2) &= 2\times 4+10 &&\justif{$(-2)^{2}=4$ ; $-5\times(-2)=+10$}\\
f(-2) &= \res{18} &&
\end{calculs}
\begin{calculs}
f\!\left(\frac{1}{2}\right) &= 2\times\left(\frac{1}{2}\right)^{2}-5\times\frac{1}{2} &&\\
f\!\left(\frac{1}{2}\right) &= 2\times\frac{1}{4}-\frac{5}{2} &&\justif{$\left(\frac{1}{2}\right)^{2}=\frac{1\times 1}{2\times 2}$}\\
f\!\left(\frac{1}{2}\right) &= \frac{1}{2}-\frac{5}{2} &&\justif{$2\times\frac{1}{4}=\frac{2}{4}$, soit $\frac{1}{2}$}\\
f\!\left(\frac{1}{2}\right) &= -\frac{4}{2} &&\justif{même dénominateur : on soustrait les numérateurs}\\
f\!\left(\frac{1}{2}\right) &= \res{-2} &&
\end{calculs}
\begin{calculs}
f(-1{,}5) &= 2\times(-1{,}5)^{2}-5\times(-1{,}5) &&\\
f(-1{,}5) &= 2\times 2{,}25+7{,}5 &&\justif{$1{,}5\times 1{,}5=2{,}25$ ; produit de deux négatifs : positif}\\
f(-1{,}5) &= 4{,}5+7{,}5 &&\\
f(-1{,}5) &= \res{12} &&
\end{calculs}
\begin{calculs}
f\!\left(\sqrt{3}\right) &= 2\times\left(\sqrt{3}\right)^{2}-5\sqrt{3} &&\\
f\!\left(\sqrt{3}\right) &= 2\times 3-5\sqrt{3} &&\justif{$\left(\sqrt{3}\right)^{2}=3$ : c'est la définition de $\sqrt{3}$}\\
f\!\left(\sqrt{3}\right) &= \res{6-5\sqrt{3}} &&\justif{valeur exacte : on ne peut pas réduire davantage}
\end{calculs}
\rem{Erreur fréquente : écrire $-2^{2}$ au lieu de $(-2)^{2}$. Sans parenthèses,
$-2^{2}=-4$ : le carré ne porte que sur $2$.}
\cas{2)}
\explique{« $-1$ est-il un antécédent de $7$ ? » revient à : a-t-on $f(-1)=7$ ?}{on
calcule l'image de $-1$ et on la compare à $7$}
\begin{calculs}
f(-1) &= 2\times(-1)^{2}-5\times(-1) &&\\
f(-1) &= 2+5 &&\justif{$(-1)^{2}=1$ ; $-5\times(-1)=+5$}\\
f(-1) &= 7 &&
\end{calculs}
\explique{$f(-1)=7$ : \res{oui}, $-1$ est un antécédent de $7$ par $f$.}{conclusion en
phrase}
\begin{calculs}
f(2) &= 2\times 2^{2}-5\times 2 &&\\
f(2) &= 8-10 &&\\
f(2) &= -2 &&
\end{calculs}
\explique{L'image de $2$ est $-2$, et non $2$ : \res{non}, $2$ n'a pas pour image
$2$.}{un nombre n'a qu'une seule image}
\end{exo}

\begin{exo}[Calculer des antécédents]
\explique{Chercher les antécédents d'un nombre $k$, c'est résoudre l'équation
$f(x)=k$.}{méthode 1 : on écrit l'équation, puis on isole $x$}
\cas{a)}
\begin{calculs}
x-7 &= 2 &&\justif{antécédents de $2$ : on résout $f(x)=2$}\\
x &= 9 &&\justif{on ajoute $7$}
\end{calculs}
\explique{Le nombre $2$ a un seul antécédent par $f$ : \res{$9$}.}{contrôle :
$f(9)=9-7=2$}
\cas{b)}
\begin{calculs}
x^{2} &= 49 &&\justif{on résout $g(x)=49$}
\end{calculs}
\explique{Le nombre $49$ a deux antécédents par $g$ : \res{$-7$ et $7$}.}{deux nombres
ont pour carré $49$ : $7^{2}=49$ et $(-7)^{2}=49$}
\cas{c)}
\begin{calculs}
5x+3 &= 0 &&\justif{on résout $h(x)=0$}\\
5x &= -3 &&\justif{on retire $3$}\\
x &= -\frac{3}{5} &&\justif{on divise par $5$}
\end{calculs}
\explique{Le nombre $0$ a un seul antécédent par $h$ : \res{$-\frac{3}{5}$}.}{valeur
exacte : la fraction, pas un arrondi}
\cas{d)}
\begin{calculs}
(x-1)(x+4) &= 0 &&\justif{on résout $k(x)=0$ : un produit nul}\\
x-1 &= 0 \quad\text{ou}\quad x+4=0 &&\justif{un produit est nul si l'un de ses facteurs est nul}\\
x &= 1 \quad\text{ou}\quad x=-4 &&
\end{calculs}
\explique{Le nombre $0$ a deux antécédents par $k$ : \res{$-4$ et $1$}.}{contrôle :
$k(1)=0\times 5$ et $k(-4)=(-5)\times 0$}
\rem{On ne développe pas : la forme factorisée est faite pour résoudre « produit
$=0$ ».}
\cas{e)}
\begin{calculs}
x^{2}+3 &= 2 &&\justif{on résout $m(x)=2$}\\
x^{2} &= -1 &&\justif{on retire $3$}
\end{calculs}
\explique{Le nombre $2$ n'a \res{aucun antécédent} par $m$.}{un carré n'est jamais
négatif}
\end{exo}

\begin{exo}[Lire des images et des antécédents]
\cas{1)}
\begin{minipage}[t]{0.56\linewidth}\raggedright
\vspace{0pt}
La courbe passe par les points $(-2\,;1)$, $(0\,;-2)$, $(3\,;2)$ et $(5\,;-1)$ :
\[
g(-2)=\res{1} \qquad g(0)=\res{-2} \qquad g(3)=\res{2} \qquad g(5)=\res{-1}
\]
\gris{Pour l'image de $-2$ : on part de $-2$ sur l'axe des abscisses, on va
verticalement jusqu'à la courbe, puis horizontalement jusqu'à l'axe des ordonnées, où on
lit l'ordonnée (tracés en pointillés). Pour $0$, le point est sur l'axe des ordonnées.}
\end{minipage}\hfill
\begin{minipage}[t]{0.40\linewidth}\centering
\vspace{0pt}
\begin{tikzpicture}[x=0.42cm,y=0.42cm]
  \repere{-4}{-3}{6}{5}
  \gradx{-3,-2,-1,1,2,3,4,5} \grady{-2,-1,1,2,3,4}
  \courbeExQuatre
  \draw[lecture] (-2,0) -- (-2,1) -- (0,1);
  \draw[lecture] (3,0) -- (3,2) -- (0,2);
  \draw[lecture] (5,0) -- (5,-1) -- (0,-1);
  \path[coulProp] (-2,1) node[croix]{} (0,-2) node[croix]{} (3,2) node[croix]{}
    (5,-1) node[croix]{};
\end{tikzpicture}
\end{minipage}
\cas{2)}
\begin{minipage}[t]{0.56\linewidth}\raggedright
\vspace{0pt}
La droite d'équation $y=1$ coupe la courbe en trois points : $1$ a trois
antécédents, \res{$-2$ ; $2$ et $4$}.\par
La droite d'équation $y=-1$ coupe la courbe en trois points : $-1$ a trois
antécédents, \res{$-1$ ; $1$ et $5$}.\par
La droite d'équation $y=-2$ touche la courbe en un seul point, $(0\,;-2)$, le plus bas :
$-2$ a un seul antécédent, \res{$0$}.\par
La droite d'équation $y=5$ ne rencontre pas la courbe : $5$ n'a \res{aucun antécédent}.\par
\gris{On trace la droite horizontale à la hauteur du nombre, puis on lit les abscisses de
\emph{tous} les points communs. La courbe ne monte pas au-dessus de $4$.}
\end{minipage}\hfill
\begin{minipage}[t]{0.40\linewidth}\centering
\vspace{0pt}
\begin{tikzpicture}[x=0.42cm,y=0.42cm]
  \repere{-4}{-3}{6}{5}
  \gradx{-3,-2,-1,1,2,3,4,5} \grady{-2,-1,1,2,3,4}
  \courbeExQuatre
  \draw[coulCourbeC,thick] (-3.6,1) -- (5.6,1) node[right,font=\footnotesize]{$y=1$};
  \draw[coulCourbeB,thick] (-3.6,-1) -- (5.6,-1) node[right,font=\footnotesize]{$y=-1$};
  \draw[coulProp,thick] (-3.6,5) -- (5.6,5) node[right,font=\footnotesize]{$y=5$};
  \draw[lecture] (-2,1) -- (-2,0) (2,1) -- (2,0) (4,1) -- (4,0)
    (-1,-1) -- (-1,0) (1,-1) -- (1,0) (5,-1) -- (5,0);
  \path[coulCourbeC] (-2,1) node[croix]{} (2,1) node[croix]{} (4,1) node[croix]{};
  \path[coulCourbeB] (-1,-1) node[croix]{} (1,-1) node[croix]{} (5,-1) node[croix]{};
  \path[coulProp] (0,-2) node[croix]{} (-3,4) node[croix]{};
\end{tikzpicture}
\end{minipage}
\cas{3)}
Les solutions de $g(x)=0$ sont les antécédents de $0$ : les abscisses des
points où la courbe coupe l'axe des abscisses. Il y en a trois (entre $-2$ et $-1$, entre
$1$ et $2$, entre $4$ et $5$) : l'équation a \res{trois solutions}.\par
\gris{On demande seulement leur nombre : leurs valeurs exactes ne se lisent pas.}
\cas{4)}
On cherche un antécédent de $4$ : la droite $y=4$ touche la courbe en
$(-3\,;4)$, donc $g(\res{-3})=4$.
\rem{Une image se lit sur l'axe des \textbf{ordonnées} ; un antécédent sur l'axe des
\textbf{abscisses}. Erreur fréquente : s'arrêter au premier point trouvé.}
\end{exo}

\begin{exo}[Points d'une courbe]
\explique{$h(x)=x^{2}+2x-8$. Un point $M(a\,;b)$ est sur $\mathcal{C}_h$ lorsque
$b=h(a)$.}{définition 1 : l'ordonnée doit être l'image de l'abscisse}
\cas{1)}
\begin{calculs}
h(1) &= 1^{2}+2\times 1-8 &&\justif{on calcule l'image de l'abscisse de $A$}\\
h(1) &= 1+2-8 &&\\
h(1) &= -5 &&
\end{calculs}
\explique{$h(1)=-5$, l'ordonnée de $A$ : \res{$A$ est sur $\mathcal{C}_h$}.}{on compare
à $y_{A}=-5$}
\begin{calculs}
h(-2) &= (-2)^{2}+2\times(-2)-8 &&\justif{abscisse de $B$, entre parenthèses}\\
h(-2) &= 4-4-8 &&\\
h(-2) &= -8 &&
\end{calculs}
\explique{$h(-2)=-8$, et non $0$ : \res{$B$ n'est pas sur $\mathcal{C}_h$}.}{le point de
$\mathcal{C}_h$ d'abscisse $-2$ est $(-2\,;-8)$}
\cas{2)}
\explique{Un point de l'axe des ordonnées a pour abscisse $0$ : on calcule $h(0)$.}{son
ordonnée est l'image de $0$}
\begin{calculs}
h(0) &= 0^{2}+2\times 0-8 &&\\
h(0) &= -8 &&
\end{calculs}
\explique{$\mathcal{C}_h$ coupe l'axe des ordonnées au point \res{$(0\,;-8)$}.}{on donne
les deux coordonnées}
\cas{3)}
\begin{calculs}
(x-2)(x+4) &= x^{2}+4x-2x-8 &&\justif{double distributivité}\\
(x-2)(x+4) &= x^{2}+2x-8 &&\justif{on reconnaît $h(x)$}
\end{calculs}
\explique{Donc, pour tout $x$, $h(x)=(x-2)(x+4)$.}{deux écritures de la même
fonction}
\explique{Un point de l'axe des abscisses a pour ordonnée $0$ : on résout
$h(x)=0$, avec la forme factorisée.}{la forme développée ne permet pas de résoudre}
\begin{calculs}
(x-2)(x+4) &= 0 &&\\
x-2 &= 0 \quad\text{ou}\quad x+4=0 &&\justif{un produit est nul si l'un de ses facteurs est nul}\\
x &= 2 \quad\text{ou}\quad x=-4 &&
\end{calculs}
\explique{$\mathcal{C}_h$ coupe l'axe des abscisses aux points \res{$(2\,;0)$ et
$(-4\,;0)$}.}{contrôle : $h(2)=4+4-8$, soit $0$}
\rem{« Le point est sur la courbe » se démontre par un calcul : ses coordonnées vérifient
$y=h(x)$. Développer sert à reconnaître $h(x)$ ; factoriser sert à résoudre $h(x)=0$.}
\end{exo}

\partie{Ensemble de définition}

\begin{exo}[Un nombre sans image]
\explique{Deux calculs sont impossibles : diviser par $0$, prendre la racine carrée d'un
nombre négatif.}{propriété 1 : on cherche lequel des quatre nombres y mène}
\cas{$f(x)=\frac{3}{x}$}
\explique{Pour $x=0$, le dénominateur est nul : \res{$0$} n'a pas d'image par $f$.}{les
autres nombres donnent $-\frac{3}{2}$ ; $\frac{3}{5}$ ; $\frac{3}{10}$}
\cas{$g(x)=\frac{x+3}{x-5}$}
\explique{$x-5=0$ pour $x=5$ : \res{$5$} n'a pas d'image par $g$.}{seul le
dénominateur compte ; le numérateur peut être nul}
\cas{$h(x)=\sqrt{x+1}$}
\explique{Pour $x=-2$, $x+1=-1$ : $\sqrt{-1}$ n'existe pas, \res{$-2$} n'a pas d'image
par $h$.}{pour $0$ ; $5$ ; $10$ : $\sqrt{1}$ ; $\sqrt{6}$ ; $\sqrt{11}$ existent}
\end{exo}

\begin{exo}[Ensembles de définition]
\explique{Quotient : on retire les valeurs qui annulent le dénominateur. Racine carrée :
le nombre sous la racine doit être positif ou nul.}{méthodes 3 et 4}
\cas{a)}
\begin{calculs}
x+6 &= 0 &&\justif{on cherche où le dénominateur s'annule}\\
x &= -6 &&\justif{$-6$ est la valeur interdite}
\end{calculs}
\explique{$f$ est définie sur \res{$\R\setminus\ens{-6}$}.}{on retire la valeur interdite
de $\R$}
\cas{b)}
\begin{calculs}
4x-12 &= 0 &&\justif{dénominateur de $g$}\\
4x &= 12 &&\justif{on ajoute $12$}\\
x &= 3 &&\justif{on divise par $4$}
\end{calculs}
\explique{$g$ est définie sur \res{$\R\setminus\ens{3}$}.}{}
\cas{c)}
\explique{Pour tout $x$, $x^{2}\ge 0$, donc $x^{2}+5\ge 5$ : le dénominateur ne s'annule
jamais. $h$ est définie sur \res{$\R$}.}{un carré est positif}
\rem{Le numérateur $x$ s'annule en $0$, mais $h(0)=\frac{0}{5}=0$ existe : seul le
dénominateur compte.}
\cas{d)}
\begin{calculs}
x^{2}-16 &= 0 &&\justif{dénominateur de $k$}\\
x^{2} &= 16 &&\justif{on ajoute $16$}
\end{calculs}
\explique{$x=-4$ ou $x=4$ : $k$ est définie sur \res{$\R\setminus\ens{-4\,;4}$}.}{deux
nombres ont pour carré $16$ : deux valeurs interdites}
\cas{e)}
\begin{calculs}
12-3x &\ge 0 &&\justif{nombre sous la racine positif ou nul}\\
-3x &\ge -12 &&\justif{on retire $12$}\\
x &\le 4 &&\justif{on divise par $-3$, négatif : le sens change}
\end{calculs}
\explique{$m$ est définie sur \res{$\intof{-\infty}{4}$}.}{contrôle : $m(4)=\sqrt{0}$
existe, $m(5)=\sqrt{-3}$ n'existe pas}
\cas{f)}
\explique{Deux conditions : $x-1\ge 0$ pour la racine, et $\sqrt{x-1}\ne 0$ pour le
dénominateur, c'est-à-dire $x-1\ne 0$.}{un quotient \emph{et} une racine : on écrit
les deux contraintes}
\begin{calculs}
x-1 &> 0 &&\justif{positif ou nul, et non nul : strictement positif}\\
x &> 1 &&
\end{calculs}
\explique{$p$ est définie sur \res{$\intoo{1}{+\infty}$}.}{borne $1$ exclue :
$p(1)=\frac{1}{\sqrt{0}}=\frac{1}{0}$ est impossible}
\end{exo}

\begin{exo}[Avec une valeur interdite]
\cas{1)}
\begin{calculs}
x-1 &= 0 &&\justif{on cherche où le dénominateur de $f(x)=\frac{6}{x-1}$ s'annule}\\
x &= 1 &&\justif{valeur interdite}
\end{calculs}
\explique{$f$ est définie sur \res{$\R\setminus\ens{1}$}.}{}
\cas{2)}
\begin{calculs}
f(0) &= \frac{6}{0-1} &&\justif{$0\ne 1$ : $0$ a une image}\\
f(0) &= \res{-6} &&\\
f(4) &= \frac{6}{4-1} &&\\
f(4) &= \frac{6}{3}=\res{2} &&\\
f(-2) &= \frac{6}{-2-1} &&\\
f(-2) &= \frac{6}{-3}=\res{-2} &&\justif{un positif divisé par un négatif : négatif}
\end{calculs}
\cas{3)}
\explique{On résout $f(x)=k$, pour $x\ne 1$.}{méthode 3}
\begin{calculs}
\frac{6}{x-1} &= 3 &&\justif{antécédents de $3$}\\
x-1 &= 2 &&\justif{$\frac{6}{\ldots}$ vaut $3$ quand le dénominateur vaut $2$}\\
x &= 3 &&\justif{$3\ne 1$ : la solution convient}
\end{calculs}
\explique{Le nombre $3$ a un antécédent : \res{$3$}.}{contrôle :
$f(3)=\frac{6}{2}$, soit $3$}
\begin{calculs}
\frac{6}{x-1} &= -6 &&\justif{antécédents de $-6$}\\
x-1 &= -1 &&\justif{$\frac{6}{-1}=-6$}\\
x &= 0 &&
\end{calculs}
\explique{Le nombre $-6$ a un antécédent : \res{$0$}.}{on le savait : $f(0)=-6$ au 2)}
\explique{Le numérateur vaut $6$, jamais $0$ : le nombre $0$ n'a \res{aucun
antécédent}.}{un quotient est nul seulement si son numérateur est nul}
\end{exo}

\begin{exo}[Avec une racine carrée]
\cas{1)}
\begin{calculs}
4x+8 &\ge 0 &&\justif{nombre sous la racine de $g(x)=\sqrt{4x+8}$ positif ou nul}\\
4x &\ge -8 &&\justif{on retire $8$}\\
x &\ge -2 &&\justif{on divise par $4$, positif : le sens ne change pas}
\end{calculs}
\explique{$g$ est définie sur \res{$\intfo{-2}{+\infty}$}.}{borne comprise :
$\sqrt{0}=0$ existe}
\cas{2)}
\begin{calculs}
g(2) &= \sqrt{4\times 2+8} &&\\
g(2) &= \sqrt{16}=\res{4} &&\\
g(-1) &= \sqrt{4\times(-1)+8} &&\\
g(-1) &= \sqrt{4}=\res{2} &&
\end{calculs}
\explique{$-3<-2$ : $-3$ n'est pas dans l'ensemble de définition,
$g(-3)=\sqrt{-4}$ \res{n'existe pas}.}{« si possible » : on vérifie d'abord que le
nombre est dans l'ensemble de définition}
\cas{3)}
\begin{calculs}
\sqrt{4x+8} &= 6 &&\justif{antécédents de $6$}\\
4x+8 &= 36 &&\justif{la racine vaut $6$ quand le nombre sous la racine vaut $36$}\\
4x &= 28 &&\justif{on retire $8$}\\
x &= 7 &&\justif{$7\ge -2$ : la solution convient}
\end{calculs}
\explique{Le nombre $6$ a un antécédent : \res{$7$}.}{contrôle : $g(7)=\sqrt{36}$, soit
$6$}
\begin{calculs}
\sqrt{4x+8} &= 0 &&\justif{antécédents de $0$}\\
4x+8 &= 0 &&\justif{seul $0$ a pour racine $0$}\\
x &= -2 &&
\end{calculs}
\explique{Le nombre $0$ a un antécédent : \res{$-2$}.}{la borne de l'ensemble de
définition}
\explique{Une racine carrée est positive : $-1$ n'a \res{aucun antécédent}.}{inutile de
calculer}
\end{exo}

\partie{Comparaison de deux fonctions}

\begin{exo}[Comparer sur un graphique]
\cas{a)}
\begin{minipage}[t]{0.56\linewidth}\raggedright
\vspace{0pt}
Les courbes se coupent en trois points, $(-2\,;1)$, $(1\,;0)$ et $(4\,;2)$ :
\[
S=\res{\ens{-2\,;1\,;4}}
\]
\gris{Les solutions de $f(x)=g(x)$ sont les \emph{abscisses} des points d'intersection.
Erreur fréquente : répondre par les ordonnées $1$ ; $0$ ; $2$.}
\end{minipage}\hfill
\begin{minipage}[t]{0.40\linewidth}\centering
\vspace{0pt}
\begin{tikzpicture}[x=0.42cm,y=0.42cm]
  \repere{-4}{-3}{6}{4}
  \gradx{-3,-2,-1,1,2,3,4,5} \grady{-2,-1,1,2,3}
  \courbesExDix
  \draw[lecture] (-2,1) -- (-2,0) (4,2) -- (4,0);
  \path[coulProp] (-2,1) node[croix]{} (1,0) node[croix]{} (4,2) node[croix]{};
\end{tikzpicture}
\end{minipage}
\cas{b)}
\begin{minipage}[t]{0.56\linewidth}\raggedright
\vspace{0pt}
$\Cf$ est strictement au-dessus de $\Cg$ de $-3$ à $-2$, puis de $1$ à $4$ :
\[
S=\res{\intfo{-3}{-2}\cup\intoo{1}{4}}
\]
\gris{Inégalité stricte : les abscisses des points d'intersection sont exclues. Le
nombre $-3$ est compris : les fonctions y sont définies et $f(-3)=3$ dépasse
$g(-3)=-1$.}
\end{minipage}\hfill
\begin{minipage}[t]{0.40\linewidth}\centering
\vspace{0pt}
\begin{tikzpicture}[x=0.42cm,y=0.42cm]
  \repere{-4}{-3}{6}{4}
  \gradx{-3,-2,-1,1,2,3,4,5} \grady{-2,-1,1,2,3}
  \courbesExDix
  \draw[inter] (-3,0) -- (-2,0);
  \draw[interB] (1,0) -- (4,0);
\end{tikzpicture}
\end{minipage}
\cas{c)}
\begin{minipage}[t]{0.56\linewidth}\raggedright
\vspace{0pt}
On trace la droite d'équation $y=2$ et on cherche où $\Cg$ est sur cette droite ou
au-dessus. $\Cg$ touche la droite en $(-1\,;2)$, son point le plus haut sur
$\intff{-3}{3}$, puis est au-dessus de la droite de $4$ à $5$ :
\[
S=\res{\ens{-1}\cup\intff{4}{5}}
\]
\gris{Seule $\Cg$ compte ici. Erreur fréquente : oublier le nombre isolé $-1$.}
\end{minipage}\hfill
\begin{minipage}[t]{0.40\linewidth}\centering
\vspace{0pt}
\begin{tikzpicture}[x=0.42cm,y=0.42cm]
  \repere{-4}{-3}{6}{4}
  \gradx{-3,-2,-1,1,2,3,4,5} \grady{-2,-1,1,2,3}
  \courbesExDix
  \draw[coulProp,thick] (-3.6,2) -- (5.6,2);
  \draw[lecture] (-1,2) -- (-1,0);
  \path[red] (-1,0) node[croix]{};
  \path[coulProp] (-1,2) node[croix]{} (4,2) node[croix]{};
  \draw[interB] (4,0) -- (5,0);
\end{tikzpicture}
\end{minipage}
\end{exo}

\begin{exo}[Signe d'une différence]
\explique{Pour chaque couple : on calcule $f(x)-g(x)$, on l'écrit sous une forme dont on
connaît le signe, puis on conclut.}{méthode 6 : signe $+$, $\Cf$ au-dessus ; signe $-$,
$\Cf$ au-dessous}
\cas{a)}
\begin{minipage}[t]{0.58\linewidth}\raggedright
\vspace{0pt}
\begin{calculs}
f(x)-g(x) &= 3x-4-(x+2) &&\justif{parenthèses autour de $g(x)$}\\
f(x)-g(x) &= 3x-4-x-2 &&\justif{le « $-$ » change les deux signes}\\
f(x)-g(x) &= 2x-6 &&
\end{calculs}
\begin{calculs}
2x-6 &> 0 &&\justif{où la différence est-elle positive ?}\\
2x &> 6 &&\\
x &> 3 &&
\end{calculs}
$f(x)-g(x)>0$ pour $x>3$, $f(x)-g(x)<0$ pour $x<3$ :\par
$\Cf$ est \res{au-dessous} de $\Cg$ sur $\intoo{-\infty}{3}$, \res{au-dessus} sur
$\intoo{3}{+\infty}$ ; les droites se coupent en \res{$(3\,;5)$}.\par
\gris{Ordonnée du point commun : $g(3)=3+2$, soit $5$ ; contrôle avec $f$ :
$f(3)=9-4$, soit $5$ aussi.}
\end{minipage}\hfill
\begin{minipage}[t]{0.38\linewidth}\centering
\vspace{0pt}
\begin{tikzpicture}[x=0.26cm,y=0.26cm]
  \repere{-1}{-3}{6}{11}
  \gradx{3} \grady{5,10}
  \draw[coulCourbe,line width=1.1pt] (1/3,-3) -- (5,11);
  \draw[coulCourbeC,line width=1.1pt] (-1,1) -- (6,8);
  \draw[lecture] (3,0) -- (3,5);
  \path[coulProp] (3,5) node[croix]{};
  \node[coulCourbe,font=\footnotesize] at (3.6,10.6) {$\Cf$};
  \node[coulCourbeC,font=\footnotesize] at (6.6,7) {$\Cg$};
\end{tikzpicture}
\end{minipage}
\cas{b)}
\begin{minipage}[t]{0.58\linewidth}\raggedright
\vspace{0pt}
\begin{calculs}
f(x)-g(x) &= x^{2}-(2x-1) &&\justif{parenthèses autour de $g(x)$}\\
f(x)-g(x) &= x^{2}-2x+1 &&\\
f(x)-g(x) &= (x-1)^{2} &&\justif{identité remarquable $a^{2}-2ab+b^{2}$}
\end{calculs}
Un carré est positif : $f(x)-g(x)\ge 0$ pour tout $x$, et $f(x)-g(x)=0$ seulement pour
$x=1$.\par
$\Cf$ est \res{au-dessus} de $\Cg$ sur $\R$, avec un seul point commun,
\res{$(1\,;1)$}.\par
\gris{$g(1)=2-1$, soit $1$. La droite touche la parabole sans la traverser.}
\end{minipage}\hfill
\begin{minipage}[t]{0.38\linewidth}\centering
\vspace{0pt}
\begin{tikzpicture}[x=0.26cm,y=0.26cm]
  \repere{-3}{-3}{4}{9}
  \gradx{-2,1,2} \grady{1,5}
  \draw[coulCourbe,line width=1.1pt,domain=-3:3,samples=50] plot (\x,{\x*\x});
  \draw[coulCourbeC,line width=1.1pt] (-1,-3) -- (4,7);
  \path[coulProp] (1,1) node[croix]{};
  \node[coulCourbe,font=\footnotesize] at (-3.6,8) {$\Cf$};
  \node[coulCourbeC,font=\footnotesize] at (4.6,6) {$\Cg$};
\end{tikzpicture}
\end{minipage}
\cas{c)}
\begin{minipage}[t]{0.58\linewidth}\raggedright
\vspace{0pt}
\begin{calculs}
f(x)-g(x) &= x^{2}-x-(-x+4) &&\justif{parenthèses autour de $g(x)$}\\
f(x)-g(x) &= x^{2}-x+x-4 &&\\
f(x)-g(x) &= x^{2}-4 &&\\
f(x)-g(x) &= (x-2)(x+2) &&\justif{$a^{2}-b^{2}=(a-b)(a+b)$}
\end{calculs}
$x+2=0$ pour $x=-2$, et $x-2=0$ pour $x=2$.\par
\gris{Le signe d'un produit se lit dans un tableau de signes : un facteur affine
$x+k$ est négatif avant sa racine, positif après.}
\end{minipage}\hfill
\begin{minipage}[t]{0.38\linewidth}\centering
\vspace{0pt}
\begin{tikzpicture}[x=0.26cm,y=0.26cm]
  \repere{-4}{-1}{4}{12}
  \gradx{-2,2} \grady{2,6,10}
  \draw[coulCourbe,line width=1.1pt,domain=-3:3.6,samples=50] plot (\x,{\x*\x-\x});
  \draw[coulCourbeC,line width=1.1pt] (-4,8) -- (4,0);
  \draw[lecture] (-2,0) -- (-2,6) (2,0) -- (2,2);
  \path[coulProp] (-2,6) node[croix]{} (2,2) node[croix]{};
  \node[coulCourbe,font=\footnotesize] at (-3.6,10.5) {$\Cf$};
  \node[coulCourbeC,font=\footnotesize] at (-3.4,6.4) {$\Cg$};
\end{tikzpicture}
\end{minipage}
\begin{center}
\begin{tikzpicture}[scale=0.75]
  \tkzTabInit[lgt=2.8,espcl=1.8,deltacl=0.6]{$x$/1,$x+2$/1,$x-2$/1,$f(x)-g(x)$/1}%
    {$-\infty$,$-2$,$2$,$+\infty$}
  \tkzTabLine{,-,z,+,t,+,}
  \tkzTabLine{,-,t,-,z,+,}
  \tkzTabLine{,+,z,-,z,+,}
\end{tikzpicture}
\end{center}
\explique{$\Cf$ est \res{au-dessous} de $\Cg$ sur $\intoo{-2}{2}$, \res{au-dessus} sur
$\intoo{-\infty}{-2}$ et sur $\intoo{2}{+\infty}$ ; elles se coupent en
\res{$(-2\,;6)$ et $(2\,;2)$}.}{signe $-$ : au-dessous ; ordonnées : $g(-2)=2+4$ et
$g(2)=-2+4$}
\rem{L'ordonnée d'un point d'intersection se calcule avec $f$ ou avec $g$ : on doit
trouver le même nombre, ce qui contrôle le calcul ($f(-2)=4+2$, soit $6$).}
\end{exo}

\begin{exo}[Une parabole et une droite]
\cas{1)}
\begin{minipage}[t]{0.58\linewidth}\raggedright
\vspace{0pt}
On calcule chaque image, par exemple :
\begin{calculs}
f(-2) &= (-2)^{2}-2\times(-2)-2 &&\\
f(-2) &= 4+4-2 &&\justif{$-2\times(-2)=+4$}\\
f(-2) &= 6 &&\\
g(-2) &= -2+2 = 0 &&
\end{calculs}
{\renewcommand{\arraystretch}{1.3}\setlength{\tabcolsep}{4pt}
\begin{tabular}{|c|*{8}{c|}}
\hline
$x$ & $-2$ & $-1$ & $0$ & $1$ & $2$ & $3$ & $4$ & $5$\\
\hline
$f(x)$ & $6$ & $1$ & $-2$ & $-3$ & $-2$ & $1$ & $6$ & $13$\\
\hline
$g(x)$ & $0$ & $1$ & $2$ & $3$ & $4$ & $5$ & $6$ & $7$\\
\hline
\end{tabular}\par}
\medskip
On place les points du tableau, puis on les relie : $\Cf$ est une parabole, $\Cg$ une
droite (figure ci-contre).\par
\gris{Le tableau de $f$ est symétrique autour de $x=1$ : bon contrôle des calculs.}
\end{minipage}\hfill
\begin{minipage}[t]{0.38\linewidth}\centering
\vspace{0pt}
\begin{tikzpicture}[x=0.32cm,y=0.32cm]
  \repere{-3}{-4}{6}{14}
  \gradx{-1,1,4} \grady{2,6,10}
  \draw[coulCourbe,line width=1.1pt,domain=-2.4:5.1,samples=60] plot (\x,{\x*\x-2*\x-2});
  \draw[coulCourbeC,line width=1.1pt] (-3,-1) -- (6,8);
  \draw[lecture] (-1,0) -- (-1,1) (4,0) -- (4,6);
  \draw[inter] (-1,0) -- (4,0);
  \path[coulCourbe] (-2,6) node[croix]{} (-1,1) node[croix]{} (0,-2) node[croix]{}
    (1,-3) node[croix]{} (2,-2) node[croix]{} (3,1) node[croix]{} (4,6) node[croix]{}
    (5,13) node[croix]{};
  \node[coulCourbe,font=\footnotesize] at (6,12.5) {$\Cf$};
  \node[coulCourbeC,font=\footnotesize] at (6.5,7) {$\Cg$};
\end{tikzpicture}
\end{minipage}
\cas{2)}
\explique{Les courbes se coupent en $(-1\,;1)$ et $(4\,;6)$ : \res{$S=\ens{-1\,;4}$}.}{on
le voit aussi dans le tableau : colonnes où $f(x)$ et $g(x)$ sont égaux}
\explique{$\Cf$ est sous $\Cg$, ou la coupe, de $-1$ à $4$ : \res{$S=\intff{-1}{4}$} (en
rouge sur l'axe).}{inégalité large : bornes comprises}
\cas{3)}
\explique{On calcule $f(x)-g(x)$, puis on développe la forme proposée : si les deux
résultats sont égaux, l'égalité est démontrée.}{on part des deux côtés}
\begin{calculs}
f(x)-g(x) &= x^{2}-2x-2-(x+2) &&\justif{parenthèses autour de $g(x)$}\\
f(x)-g(x) &= x^{2}-2x-2-x-2 &&\\
f(x)-g(x) &= x^{2}-3x-4 &&
\end{calculs}
\begin{calculs}
(x+1)(x-4) &= x^{2}-4x+x-4 &&\justif{double distributivité}\\
(x+1)(x-4) &= x^{2}-3x-4 &&
\end{calculs}
\explique{Les deux expressions sont égales : \res{$f(x)-g(x)=(x+1)(x-4)$}.}{}
\cas{4)}
\explique{$x+1=0$ pour $x=-1$ et $x-4=0$ pour $x=4$.}{on place $-1$ et $4$ dans l'ordre
sur la première ligne}
\begin{center}
\begin{tikzpicture}[scale=0.8]
  \tkzTabInit[lgt=3,espcl=2,deltacl=0.6]{$x$/1,$x+1$/1,$x-4$/1,$f(x)-g(x)$/1}%
    {$-\infty$,$-1$,$4$,$+\infty$}
  \tkzTabLine{,-,z,+,t,+,}
  \tkzTabLine{,-,t,-,z,+,}
  \tkzTabLine{,+,z,-,z,+,}
\end{tikzpicture}
\end{center}
\explique{Dernière ligne : règle des signes d'un produit.}{$(-)\times(-)$ donne $+$ ;
$(+)\times(-)$ donne $-$}
\cas{5)}
\explique{$\Cf$ est \res{au-dessus} de $\Cg$ sur $\intoo{-\infty}{-1}$ et sur
$\intoo{4}{+\infty}$, \res{au-dessous} sur $\intoo{-1}{4}$ ; elles se coupent en
$(-1\,;1)$ et $(4\,;6)$.}{signe $+$ : au-dessus ; signe $-$ : au-dessous}
\explique{C'est ce que la figure du 2) laissait voir : le calcul le
\emph{démontre}.}{le graphique conjecture, le calcul démontre}
\end{exo}

\begin{exo}[Une fonction quotient]
\cas{1)}
\begin{calculs}
x-3 &= 0 &&\justif{dénominateur de $f(x)=\frac{2x+1}{x-3}$}\\
x &= 3 &&\justif{valeur interdite}
\end{calculs}
\explique{\res{$D=\R\setminus\ens{3}$}.}{}
\cas{2)}
\begin{calculs}
f(0) &= \frac{2\times 0+1}{0-3} &&\\
f(0) &= \res{-\frac{1}{3}} &&\justif{valeur exacte}\\
f(4) &= \frac{2\times 4+1}{4-3} &&\\
f(4) &= \frac{9}{1}=\res{9} &&\\
f(2) &= \frac{2\times 2+1}{2-3} &&\\
f(2) &= \frac{5}{-1}=\res{-5} &&
\end{calculs}
\cas{3)}
\explique{Pour chaque nombre $k$, on résout $f(x)=k$ avec $x\ne 3$.}{méthode 3}
\begin{calculs}
\frac{2x+1}{x-3} &= 1 &&\justif{antécédents de $1$}\\
2x+1 &= x-3 &&\justif{un quotient vaut $1$ quand numérateur et dénominateur sont égaux}\\
x &= -4 &&\justif{on retire $x$, puis $1$ ; $-4\ne 3$}
\end{calculs}
\explique{$1$ a un antécédent : \res{$-4$}.}{contrôle : $f(-4)=\frac{-7}{-7}$, soit $1$}
\begin{calculs}
\frac{2x+1}{x-3} &= 0 &&\justif{antécédents de $0$}\\
2x+1 &= 0 &&\justif{un quotient est nul quand son numérateur est nul}\\
x &= -\frac{1}{2} &&
\end{calculs}
\explique{$0$ a un antécédent : \res{$-\frac{1}{2}$}.}{$-\frac{1}{2}\ne 3$ : la solution
convient}
\begin{calculs}
\frac{2x+1}{x-3} &= 2 &&\justif{antécédents de $2$}\\
2x+1 &= 2(x-3) &&\justif{on multiplie par $x-3$, non nul}\\
2x+1 &= 2x-6 &&\\
1 &= -6 &&\justif{on retire $2x$ : égalité fausse}
\end{calculs}
\explique{L'équation n'a pas de solution : $2$ n'a \res{aucun antécédent}.}{les $x$
disparaissent et il reste une égalité fausse}
\cas{4)}
\explique{On met $2+\frac{7}{x-3}$ au même dénominateur $x-3$.}{on part de l'expression
proposée et on retrouve $f(x)$}
\begin{calculs}
2+\frac{7}{x-3} &= \frac{2\fois{(x-3)}}{1\fois{(x-3)}}+\frac{7}{x-3} &&\justif{$2=\frac{2}{1}$}\\
2+\frac{7}{x-3} &= \frac{2x-6+7}{x-3} &&\justif{on additionne les numérateurs}\\
2+\frac{7}{x-3} &= \frac{2x+1}{x-3} &&
\end{calculs}
\explique{Pour tout $x$ de $D$, \res{$f(x)=2+\frac{7}{x-3}$}.}{}
\cas{5)}
\explique{Comparer $\Cf$ à la droite d'équation $y=2$, c'est étudier le signe de
$f(x)-2$.}{méthode 6, avec $g(x)=2$}
\begin{calculs}
f(x)-2 &= \frac{7}{x-3} &&\justif{d'après le 4)}
\end{calculs}
\explique{Le numérateur $7$ est positif : $f(x)-2$ a le signe de $x-3$.}{un quotient de
deux positifs est positif ; d'un positif par un négatif, négatif}
\begin{itemize}
  \item pour $x>3$, $f(x)-2>0$ : $\Cf$ est \res{au-dessus} de la droite $y=2$ ;
  \item pour $x<3$, $f(x)-2<0$ : $\Cf$ est \res{au-dessous} de la droite $y=2$.
\end{itemize}
\explique{$f(x)-2$ n'est jamais nul : la courbe ne coupe jamais la droite.}{on retrouve le
3) : $2$ n'a pas d'antécédent}
\cas{6)}
\begin{minipage}[t]{0.58\linewidth}\raggedright
\vspace{0pt}
\begin{calculs}
f(-4) &= \frac{-8+1}{-4-3}=\res{1} &&\\
f(-1) &= \frac{-2+1}{-1-3}=\res{0{,}25} &&\\
f(5) &= \frac{10+1}{5-3}=\res{5{,}5} &&\\
f(10) &= \frac{20+1}{10-3}=\res{3} &&
\end{calculs}
Avec le 2) : $f(0)\approx -0{,}33$ ; $f(2)=-5$ ; $f(4)=9$.\par
\gris{On place ces sept points, sans relier à travers la droite verticale $x=3$, où $f$
n'est pas définie.}\par
La courbe a deux branches, de part et d'autre de la droite $x=3$ ; elles s'approchent de
la droite $y=2$, la branche de gauche par-dessous, celle de droite par-dessus, comme l'annonçait le 5).
\end{minipage}\hfill
\begin{minipage}[t]{0.38\linewidth}\centering
\vspace{0pt}
\begin{tikzpicture}[x=0.25cm,y=0.25cm]
  \repere{-6}{-7}{12}{11}
  \gradx{-4,3,5,10} \grady{-5,2,5,9}
  \draw[dashed] (3,-7) -- (3,11);
  \draw[coulCourbeC,line width=1.1pt] (-6,2) -- (12,2);
  \draw[coulCourbe,line width=1.1pt,domain=-6:2.125,samples=60] plot (\x,{(2*\x+1)/(\x-3)});
  \draw[coulCourbe,line width=1.1pt,domain=3.875:12,samples=60] plot (\x,{(2*\x+1)/(\x-3)});
  \path[coulCourbe] (-4,1) node[croix]{} (-1,0.25) node[croix]{} (0,-1/3) node[croix]{}
    (2,-5) node[croix]{} (4,9) node[croix]{} (5,5.5) node[croix]{} (10,3) node[croix]{};
  \node[coulCourbe,font=\footnotesize] at (6,9.5) {$\Cf$};
  \node[coulCourbeC,font=\footnotesize] at (11,3.2) {$y=2$};
\end{tikzpicture}
\end{minipage}
\end{exo}

\partie{Sens de variation}

\begin{exo}[Lecture graphique complète]
\cas{1)}
\begin{minipage}[t]{0.56\linewidth}\raggedright
\vspace{0pt}
Antécédents de $0$ : \res{$-4$ ; $0$ et $4$}.\par
\gris{Points où $\Cf$ coupe l'axe des abscisses (la droite $y=0$).}\par
Antécédents de $1$ : \res{$-3$ ; $-1$ et $5$}.\par
\gris{La droite $y=1$ (en vert) coupe $\Cf$ trois fois.}\par
Antécédents de $-2$ : \res{$1$ et $3$}.\par
\gris{La droite $y=-2$ (en rouge foncé) coupe $\Cf$ deux fois.}\par
Antécédents de $4$ : \res{aucun}.\par
\gris{La droite $y=4$ ne rencontre pas $\Cf$, qui ne dépasse pas $3$.}
\end{minipage}\hfill
\begin{minipage}[t]{0.40\linewidth}\centering
\vspace{0pt}
\begin{tikzpicture}[x=0.4cm,y=0.4cm]
  \repereExQuatorze
  \courbeExQuatorze
  \draw[coulCourbeC,thick] (-4.6,1) -- (6.6,1);
  \draw[coulCourbeB,thick] (-4.6,-2) -- (6.6,-2);
  \draw[coulProp,thick] (-4.6,4) -- (6.6,4) node[right,font=\footnotesize]{$y=4$};
  \draw[lecture] (-3,1) -- (-3,0) (-1,1) -- (-1,0) (5,1) -- (5,0) (1,-2) -- (1,0)
    (3,-2) -- (3,0);
  \path[coulProp] (0,0) node[croix]{} (4,0) node[croix]{};
  \path[coulCourbeC] (-3,1) node[croix]{} (-1,1) node[croix]{} (5,1) node[croix]{};
  \path[coulCourbeB] (1,-2) node[croix]{} (3,-2) node[croix]{};
\end{tikzpicture}
\end{minipage}
\cas{2) et 3)}
\begin{minipage}[t]{0.56\linewidth}\raggedright
\vspace{0pt}
\textbf{2)} Le point le plus haut de $\Cf$ est $(6\,;3)$, le plus bas $(2\,;-3)$ :\par
le \textbf{maximum} de $f$ est \res{$3$, atteint en $6$} ;\par
le \textbf{minimum} de $f$ est \res{$-3$, atteint en $2$}.\par
\gris{Le maximum est une valeur de $f(x)$ (une ordonnée) ; « atteint en » donne
l'abscisse. Le sommet $(-2\,;2)$ n'est pas le maximum : $\Cf$ monte plus haut en $6$.}

\medskip
\textbf{3)} En suivant $\Cf$ de gauche à droite : elle monte de $-4$ à $-2$, descend de
$-2$ à $2$, puis monte de $2$ à $6$.\par
\gris{Dans le tableau : les abscisses où le sens change sur la première ligne, les
ordonnées correspondantes au bout des flèches.}
\end{minipage}\hfill
\begin{minipage}[t]{0.40\linewidth}\centering
\vspace{0pt}
\begin{tikzpicture}[x=0.4cm,y=0.4cm]
  \repereExQuatorze
  \courbeExQuatorze
  \draw[lecture] (6,0) -- (6,3) -- (0,3) (2,0) -- (2,-3) -- (0,-3) (-2,0) -- (-2,2);
  \path[coulProp] (6,3) node[croix]{} (2,-3) node[croix]{} (-2,2) node[croix]{};
\end{tikzpicture}
\end{minipage}
\begin{center}
\begin{tikzpicture}[scale=0.8]
  \tkzTabInit[lgt=1.6,espcl=2]{$x$/1,$f(x)$/2}{$-4$,$-2$,$2$,$6$}
  \tkzTabVar{-/$0$,+/$2$,-/$-3$,+/$3$}
\end{tikzpicture}
\end{center}
\explique{$f$ est \res{croissante} sur $\intff{-4}{-2}$, \res{décroissante} sur
$\intff{-2}{2}$, \res{croissante} sur $\intff{2}{6}$.}{la même chose en phrases}
\cas{4)}
\begin{minipage}[t]{0.56\linewidth}\raggedright
\vspace{0pt}
$f(x)\ge 0$ : $\Cf$ est sur l'axe des abscisses ou au-dessus, de $-4$ à $0$ et de $4$ à
$6$ :
\[
S=\res{\intff{-4}{0}\cup\intff{4}{6}}
\]
\gris{Inégalité large : $-4$ ; $0$ et $4$, où $f(x)=0$, sont compris ; $6$ est la
borne de l'ensemble de définition.}
\end{minipage}\hfill
\begin{minipage}[t]{0.40\linewidth}\centering
\vspace{0pt}
\begin{tikzpicture}[x=0.4cm,y=0.4cm]
  \repereExQuatorze
  \courbeExQuatorze
  \draw[inter] (-4,0) -- (0,0);
  \draw[interB] (4,0) -- (6,0);
\end{tikzpicture}
\end{minipage}

\medskip
\begin{minipage}[t]{0.56\linewidth}\raggedright
\vspace{0pt}
$f(x)<1$ : $\Cf$ est strictement sous la droite $y=1$, de $-4$ à $-3$ et de $-1$ à $5$ :
\[
S=\res{\intfo{-4}{-3}\cup\intoo{-1}{5}}
\]
\gris{Inégalité stricte : $-3$ ; $-1$ et $5$, où $f(x)=1$, sont exclus ; $-4$ est
compris, car $f(-4)=0<1$.}
\end{minipage}\hfill
\begin{minipage}[t]{0.40\linewidth}\centering
\vspace{0pt}
\begin{tikzpicture}[x=0.4cm,y=0.4cm]
  \repereExQuatorze
  \courbeExQuatorze
  \draw[coulCourbeC,thick] (-4.6,1) -- (6.6,1) node[right,font=\footnotesize]{$y=1$};
  \draw[lecture] (-3,1) -- (-3,0) (-1,1) -- (-1,0) (5,1) -- (5,0);
  \draw[inter] (-4,0) -- (-3,0);
  \draw[interB] (-1,0) -- (5,0);
\end{tikzpicture}
\end{minipage}
\cas{5)}
\explique{Si $-2\le x\le 2$, alors \res{$-3$}$\;\le f(x)\le\;$\res{$2$}.}{tableau du
3) : sur $\intff{-2}{2}$, $f$ descend de $2$ à $-3$}
\explique{Si $2\le x\le 5$, alors \res{$-3$}$\;\le f(x)\le\;$\res{$1$}.}{sur
$\intff{2}{5}$, $f$ monte de $f(2)=-3$ à $f(5)=1$}
\rem{Erreur fréquente : écrire $-3\le f(x)\le 3$ pour $2\le x\le 5$. On s'arrête en $5$ :
la plus grande valeur est $f(5)=1$, pas le maximum $3$ atteint en $6$.}
\cas{6)}
\begin{minipage}[t]{0.56\linewidth}\raggedright
\vspace{0pt}
$f(x)=-x+1$ : on cherche les points communs à $\Cf$ et à $(D)$. Il n'y en a qu'un,
$(3\,;-2)$ :
\[
S=\res{\ens{3}}
\]
\gris{Contrôle : $-3+1=-2$, donc $(3\,;-2)$ est bien sur $(D)$.}\par
$f(x)\le -x+1$ : $\Cf$ est sous $(D)$, ou la coupe, de $-4$ à $3$ :
\[
S=\res{\intff{-4}{3}}
\]
\end{minipage}\hfill
\begin{minipage}[t]{0.40\linewidth}\centering
\vspace{0pt}
\begin{tikzpicture}[x=0.4cm,y=0.4cm]
  \repereExQuatorze
  \courbeExQuatorze
  \draw[coulCourbeC,line width=1.1pt] (-3,4) -- (5,-4);
  \node[coulCourbeC,font=\footnotesize] at (-2.2,4.5) {$(D)$};
  \draw[lecture] (3,-2) -- (3,0);
  \path[coulProp] (3,-2) node[croix]{};
  \draw[inter] (-4,0) -- (3,0);
\end{tikzpicture}
\end{minipage}
\end{exo}

\begin{exo}[Du tableau de variations à la courbe]
\begin{minipage}[t]{0.58\linewidth}\raggedright
\vspace{0pt}
\cas{1)}
\explique{L'ensemble de définition est \res{$\intff{-3}{4}$}.}{première ligne du tableau :
de la plus petite à la plus grande valeur de $x$}
\cas{2)}
\explique{$g$ est \res{croissante} sur $\intff{-3}{0}$ (de $-1$ à $3$), \res{décroissante}
sur $\intff{0}{2}$ (de $3$ à $-2$), puis \res{croissante} sur $\intff{2}{4}$ (de $-2$ à
$1$).}{flèche qui monte : croissante ; flèche qui descend : décroissante}
\cas{3)}
\explique{Le maximum est \res{$3$, atteint en $0$} ; le minimum est \res{$-2$, atteint en
$2$}.}{on compare les valeurs aux bouts des flèches : $-1$ ; $3$ ; $-2$ ; $1$}
\cas{4)}
\explique{Une courbe possible ci-contre : elle passe par $(-3\,;-1)$, $(0\,;3)$,
$(2\,;-2)$ et $(4\,;1)$, et respecte les variations.}{on place d'abord les points du
tableau, puis on les relie en montant ou en descendant, sommets arrondis}
\end{minipage}\hfill
\begin{minipage}[t]{0.38\linewidth}\centering
\vspace{0pt}
\begin{tikzpicture}[x=0.4cm,y=0.4cm]
  \repere{-4}{-3}{5}{4}
  \gradx{-3,-2,-1,1,2,3,4} \grady{-2,-1,1,2,3}
  \draw[coulCourbe,line width=1.1pt] (-3,-1)
    .. controls (-2,-1) and (-1,3) .. (0,3)
    .. controls (0.667,3) and (1.333,-2) .. (2,-2)
    .. controls (2.667,-2) and (3.333,0) .. (4,1);
  \path[coulCourbe] (-3,-1) node[croix]{} (0,3) node[croix]{} (2,-2) node[croix]{}
    (4,1) node[croix]{};
  \node[coulCourbe,font=\footnotesize] at (4.4,1.8) {$\Cg$};
\end{tikzpicture}
\end{minipage}
\rem{Un tableau de variations ne donne que l'allure de la courbe : il y a une infinité de
courbes possibles. Erreur fréquente : prendre le premier nombre de la ligne de $g(x)$,
$-1$, pour le minimum.}
\end{exo}

\begin{exo}[Le bénéfice d'un atelier]
\cas{1)}
\explique{$B(x)=-x^{2}+12x-20$ : on remplace $x$ par chaque entier de $1$ à $11$, par
exemple :}{$-x^{2}$ : on élève au carré, puis on prend l'opposé}
\begin{calculs}
B(1) &= -1^{2}+12\times 1-20 &&\justif{$-1^{2}=-1$ : le carré ne porte que sur $1$}\\
B(1) &= -1+12-20 = -9 &&\\
B(6) &= -6^{2}+12\times 6-20 &&\\
B(6) &= -36+72-20 = 16 &&
\end{calculs}
\begin{center}
{\renewcommand{\arraystretch}{1.3}\setlength{\tabcolsep}{4pt}
\begin{tabular}{|c|*{11}{c|}}
\hline
$x$ & $1$ & $2$ & $3$ & $4$ & $5$ & $6$ & $7$ & $8$ & $9$ & $10$ & $11$\\
\hline
$B(x)$ & $-9$ & $0$ & $7$ & $12$ & $15$ & $16$ & $15$ & $12$ & $7$ & $0$ & $-9$\\
\hline
\end{tabular}}
\end{center}
\rem{Le tableau est symétrique autour de $x=6$ : bon contrôle des calculs.}
\cas{2)}
\begin{minipage}[t]{0.50\linewidth}\raggedright
\vspace{0pt}
On place les onze points du tableau, puis on les relie : une parabole « tournée vers le
bas », symétrique par rapport à la droite $x=6$.\par
\gris{Unités de l'énoncé : $1$~cm pour une dizaine en abscisses, $1$~cm pour $2$
centaines d'euros en ordonnées (figure ci-contre à l'échelle $\frac{1}{2}$). Le repère
n'est pas orthonormé : les deux axes portent des grandeurs différentes.}
\end{minipage}\hfill
\begin{minipage}[t]{0.46\linewidth}\centering
\vspace{0pt}
\begin{tikzpicture}[x=0.5cm,y=0.25cm]
  \foreach \k in {-10,-8,...,18} \draw[grille] (0,\k) -- (12,\k);
  \foreach \k in {0,...,12} \draw[grille] (\k,-10) -- (\k,18);
  \draw[-{Stealth[length=2mm]},thick] (0,0) -- (12.5,0) node[below left]{$x$};
  \draw[-{Stealth[length=2mm]},thick] (0,-10) -- (0,19.5) node[below left]{$y$};
  \foreach \k in {2,4,6,8,10} \node[below,font=\scriptsize] at (\k,0) {$\k$};
  \foreach \k in {-8,-4,4,8,12,16} \node[left,font=\scriptsize] at (0,\k) {$\k$};
  \draw[coulCourbe,line width=1.1pt,domain=1:11,samples=60] plot (\x,{-\x*\x+12*\x-20});
  \draw[lecture] (6,0) -- (6,16) -- (0,16);
  \draw[inter] (2,0) -- (10,0);
  \path[coulCourbe] (1,-9) node[croix]{} (2,0) node[croix]{} (3,7) node[croix]{}
    (4,12) node[croix]{} (5,15) node[croix]{} (6,16) node[croix]{} (7,15) node[croix]{}
    (8,12) node[croix]{} (9,7) node[croix]{} (10,0) node[croix]{} (11,-9) node[croix]{};
\end{tikzpicture}
\end{minipage}
\cas{3)}
\begin{center}
\begin{tikzpicture}[scale=0.8]
  \tkzTabInit[lgt=1.6,espcl=2.4]{$x$/1,$B(x)$/2}{$1$,$6$,$11$}
  \tkzTabVar{-/$-9$,+/$16$,-/$-9$}
\end{tikzpicture}
\end{center}
\explique{$B$ est croissante sur $\intff{1}{6}$, décroissante sur $\intff{6}{11}$.}{la
courbe monte jusqu'au sommet $(6\,;16)$, puis descend}
\cas{4)}
\explique{Le maximum de $B$ est $16$, atteint en $6$ : le bénéfice maximal est \res{$16$
centaines d'euros, soit $1\,600$~€}, pour \res{$60$ trottinettes}.}{$x=6$ dizaines ;
$B=16$ centaines d'euros (pointillés)}
\cas{5)}
\explique{$B(x)\ge 0$ : la courbe est sur l'axe des abscisses ou au-dessus, de $2$ à $10$
(en rouge). L'atelier ne perd pas d'argent lorsqu'il vend \res{entre $20$ et $100$
trottinettes}.}{$B(2)=0$ et $B(10)=0$ : bornes comprises ; on traduit les dizaines en
trottinettes}
\rem{Attention aux unités : $x$ est en \emph{dizaines} de trottinettes, $B(x)$ en
\emph{centaines} d'euros. La réponse se donne dans les unités du problème.}
\end{exo}

\begin{exo}[Sens de variation par le calcul]
\explique{Dans tout l'exercice, $a<b$, donc $a-b<0$.}{méthode 7 : calculer $f(a)-f(b)$,
factoriser pour faire apparaître $a-b$, conclure sur le signe}
\cas{1)}
\begin{calculs}
f(a)-f(b) &= (4a-7)-(4b-7) &&\justif{parenthèses autour de $f(b)$}\\
f(a)-f(b) &= 4a-7-4b+7 &&\\
f(a)-f(b) &= 4a-4b &&\\
f(a)-f(b) &= 4(a-b) &&\justif{on factorise pour faire apparaître $a-b$}
\end{calculs}
\explique{$4>0$ et $a-b<0$, donc $f(a)-f(b)<0$, c'est-à-dire $f(a)<f(b)$.}{un positif
fois un négatif : négatif}
\explique{$f$ est \res{strictement croissante} sur $\R$.}{$a<b$ entraîne $f(a)<f(b)$ :
$f$ conserve l'ordre}
\cas{2)}
\begin{calculs}
g(a)-g(b) &= (-3a+2)-(-3b+2) &&\justif{parenthèses autour de $g(b)$}\\
g(a)-g(b) &= -3a+2+3b-2 &&\\
g(a)-g(b) &= -3a+3b &&\\
g(a)-g(b) &= -3(a-b) &&\justif{on fait apparaître $a-b$}
\end{calculs}
\explique{$-3<0$ et $a-b<0$, donc $g(a)-g(b)>0$, c'est-à-dire $g(a)>g(b)$.}{deux
négatifs : produit positif}
\explique{$g$ est \res{strictement décroissante} sur $\R$.}{$a<b$ entraîne $g(a)>g(b)$ :
$g$ inverse l'ordre}
\cas{3)}
\explique{Le même calcul donne $(ma+p)-(mb+p)=m(a-b)$, avec $a-b<0$.}{le signe ne dépend
plus que de $m$}
\begin{itemize}
  \item si $m>0$, $m(a-b)<0$ : la fonction est \res{strictement croissante} ;
  \item si $m<0$, $m(a-b)>0$ : elle est \res{strictement décroissante} ;
  \item si $m=0$, $m(a-b)=0$ : elle est \res{constante} (elle vaut $p$).
\end{itemize}
\cas{4)}
\begin{calculs}
u(a)-u(b) &= \left(a^{2}+4a\right)-\left(b^{2}+4b\right) &&\justif{parenthèses autour de $u(b)$}\\
u(a)-u(b) &= a^{2}-b^{2}+4a-4b &&\justif{on regroupe les carrés, puis le reste}\\
u(a)-u(b) &= (a-b)(a+b)+4(a-b) &&\justif{$a^{2}-b^{2}=(a-b)(a+b)$}\\
u(a)-u(b) &= (a-b)(a+b+4) &&\justif{facteur commun $a-b$}
\end{calculs}
\explique{Sur $\intfo{0}{+\infty}$, $a\ge 0$ et $b>0$, donc $a+b+4>0$.}{c'est là que
l'intervalle sert}
\explique{$a-b<0$ et $a+b+4>0$ : $u(a)-u(b)<0$, donc $u(a)<u(b)$. $u$ est
\res{strictement croissante} sur $\intfo{0}{+\infty}$.}{un négatif fois un positif}
\cas{5)}
\begin{calculs}
v(a)-v(b) &= \left(1-\frac{1}{a}\right)-\left(1-\frac{1}{b}\right) &&\justif{parenthèses autour de $v(b)$}\\
v(a)-v(b) &= 1-\frac{1}{a}-1+\frac{1}{b} &&\\
v(a)-v(b) &= -\frac{1\fois{b}}{a\fois{b}}+\frac{1\fois{a}}{b\fois{a}} &&\justif{même dénominateur $ab$}\\
v(a)-v(b) &= \frac{-b+a}{ab} &&\\
v(a)-v(b) &= \frac{a-b}{ab} &&
\end{calculs}
\explique{$a-b<0$ et $ab>0$ (produit de deux nombres positifs) : $v(a)-v(b)<0$, donc
$v(a)<v(b)$. $v$ est \res{strictement croissante} sur $\intoo{0}{+\infty}$.}{un négatif
divisé par un positif : négatif}
\rem{Toujours la même démarche : calculer $f(a)-f(b)$, \emph{factoriser} pour faire
apparaître $a-b$ (ou le mettre au numérateur), puis conclure sur le signe. Erreur
fréquente : oublier les parenthèses autour de $f(b)$.}
\end{exo}

\partie{Fonction carré}

\begin{exo}[Équations avec des carrés]
\explique{On se ramène à $x^{2}=k$ (ou $A^{2}=k$) : deux solutions si $k>0$, une si
$k=0$, aucune si $k<0$.}{définition 5}
\cas{a)}
\explique{$x^{2}=36$ et $36>0$ : $x=-6$ ou $x=6$, \res{$S=\ens{-6\,;6}$}.}{$6^{2}=36$ et
$(-6)^{2}=36$}
\cas{b)}
\explique{$x^{2}=11$ et $11>0$ : $x=-\sqrt{11}$ ou $x=\sqrt{11}$,
\res{$S=\ens{-\sqrt{11}\,;\sqrt{11}}$}.}{$11$ n'est pas un carré parfait : valeurs
exactes avec la racine}
\cas{c)}
\begin{calculs}
2x^{2}-5 &= x^{2}+4 &&\\
2x^{2}-x^{2} &= 4+5 &&\justif{les $x^{2}$ d'un côté, les nombres de l'autre}\\
x^{2} &= 9 &&
\end{calculs}
\explique{$9>0$ : \res{$S=\ens{-3\,;3}$}.}{contrôle : $2\times 9-5=13$ et $9+4=13$}
\cas{d)}
\explique{$(3x+1)^{2}=16$ est de la forme $A^{2}=16$ : $A=4$ ou $A=-4$.}{on ne développe
pas : on résout comme $x^{2}=16$}
\begin{calculs}
3x+1 &= 4 &&\justif{premier cas}\\
3x &= 3 &&\\
x &= 1 &&
\end{calculs}
\begin{calculs}
3x+1 &= -4 &&\justif{second cas}\\
3x &= -5 &&\\
x &= -\frac{5}{3} &&
\end{calculs}
\explique{\res{$S=\ens{-\frac{5}{3}\,;1}$}.}{contrôle : $(3\times 1+1)^{2}=4^{2}$, soit
$16$}
\cas{e)}
\begin{calculs}
x^{2}+9 &= 0 &&\\
x^{2} &= -9 &&\justif{on retire $9$}
\end{calculs}
\explique{Un carré n'est jamais négatif : \res{$S=\varnothing$}.}{$k=-9<0$ : aucune
solution}
\end{exo}

\begin{exo}[Comparer et encadrer des carrés]
\explique{La fonction carré est décroissante sur $\intof{-\infty}{0}$ (elle inverse
l'ordre) et croissante sur $\intfo{0}{+\infty}$ (elle conserve l'ordre).}{propriété 3 :
on regarde d'abord le signe des nombres}
\cas{1)}
\explique{$-4{,}2<-3{,}9\le 0$ : deux négatifs, la fonction carré inverse l'ordre, donc
\res{$(-4{,}2)^{2}>(-3{,}9)^{2}$}.}{décroissante sur $\intof{-\infty}{0}$}
\explique{$0\le 1{,}05<1{,}5$ : deux positifs, la fonction carré conserve l'ordre, donc
\res{$1{,}05^{2}<1{,}5^{2}$}.}{croissante sur $\intfo{0}{+\infty}$}
\cas{2)}
\explique{$3\le x\le 6$ donne $3^{2}\le x^{2}\le 6^{2}$, soit
\res{$9\le x^{2}\le 36$}.}{nombres positifs : l'ordre est conservé}
\explique{$-5\le x\le -2$ donne $(-2)^{2}\le x^{2}\le (-5)^{2}$, soit
\res{$4\le x^{2}\le 25$}.}{nombres négatifs : l'ordre s'inverse, les bornes changent de
place}
\explique{$-3\le x\le 1$ donne \res{$0\le x^{2}\le 9$}.}{l'intervalle contient $0$ :
$x^{2}$ descend de $9$ à $0$, puis monte jusqu'à $1$ ; le plus grand est $9$}
\rem{Piège du dernier cas : écrire $1\le x^{2}\le 9$. Dès que l'intervalle contient $0$,
la plus petite valeur de $x^{2}$ est $0$ (pour $x=0$).}
\end{exo}

\begin{exo}[Deux paraboles]
\cas{1)}
\begin{calculs}
f(x)-g(x) &= 2x^{2}+x+7-\left(x^{2}-3x+3\right) &&\justif{parenthèses autour de $g(x)$}\\
f(x)-g(x) &= 2x^{2}+x+7-x^{2}+3x-3 &&\justif{le « $-$ » change les trois signes}\\
f(x)-g(x) &= x^{2}+4x+4 &&
\end{calculs}
\begin{calculs}
(x+2)^{2} &= x^{2}+2\times x\times 2+2^{2} &&\justif{$(a+b)^{2}=a^{2}+2ab+b^{2}$}\\
(x+2)^{2} &= x^{2}+4x+4 &&
\end{calculs}
\explique{Les deux expressions sont égales : \res{$f(x)-g(x)=(x+2)^{2}$}.}{}
\cas{2)}
\explique{Un carré est positif : $f(x)-g(x)\ge 0$ pour tout $x$. \res{$\Cf$ est au-dessus
de $\Cg$ sur $\R$}.}{signe $+$ : au-dessus}
\explique{$f(x)-g(x)=0$ seulement pour $x+2=0$, soit $x=-2$.}{un carré est nul quand le
nombre élevé au carré est nul}
\begin{calculs}
f(-2) &= 2\times(-2)^{2}+(-2)+7 &&\justif{ordonnée du point commun}\\
f(-2) &= 8-2+7 &&\\
f(-2) &= 13 &&
\end{calculs}
\explique{Le point commun est \res{$(-2\,;13)$}.}{contrôle : $g(-2)=4+6+3$, soit $13$}
\end{exo}

\partie{Fonction cube}

\begin{exo}[Variations de la fonction cube]
\cas{1)}
\begin{calculs}
(a-b)\left(a^{2}+ab+b^{2}\right) &= a^{3}+a^{2}b+ab^{2}-a^{2}b-ab^{2}-b^{3} &&\justif{chaque terme de la première parenthèse multiplie la seconde}\\
(a-b)\left(a^{2}+ab+b^{2}\right) &= \res{a^{3}-b^{3}} &&\justif{les termes $a^{2}b$ et $ab^{2}$ s'annulent}
\end{calculs}
\explique{Comparer $a^{3}$ et $b^{3}$, c'est étudier le signe de $a^{3}-b^{3}$, donc celui
du produit $(a-b)\left(a^{2}+ab+b^{2}\right)$.}{même idée que pour la fonction carré, avec
$a-b<0$}
\cas{2)}
\explique{Si $a<b\le 0$ : $a^{2}>0$ (car $a<0$), $b^{2}\ge 0$, et $ab\ge 0$.}{produit de
deux nombres négatifs (ou nuls) : positif}
\explique{Une somme de nombres positifs, dont l'un strictement : $a^{2}+ab+b^{2}>0$.}{}
\explique{$a-b<0$ et $a^{2}+ab+b^{2}>0$ : $a^{3}-b^{3}$ est \res{négatif}, donc
$a^{3}<b^{3}$.}{un négatif fois un positif}
\explique{La fonction cube est \res{strictement croissante} sur
$\intof{-\infty}{0}$.}{$a<b$ entraîne $a^{3}<b^{3}$}
\cas{3)}
\explique{Si $0\le a<b$ : $a^{2}\ge 0$ ; $b^{2}>0$ et $ab\ge 0$, donc
$a^{2}+ab+b^{2}>0$ ; de nouveau $a^{3}-b^{3}<0$.}{même raisonnement : deux nombres
positifs}
\explique{La fonction cube est strictement croissante sur $\intfo{0}{+\infty}$.}{}
\explique{Enfin, si $a<0<b$ : $a^{3}<0$ et $b^{3}>0$, donc $a^{3}<b^{3}$.}{le cube d'un
négatif est négatif ; il reste à comparer deux nombres de part et d'autre de $0$}
\explique{La fonction cube est \res{strictement croissante sur $\R$}.}{elle conserve
toujours l'ordre}
\rem{C'est la différence avec la fonction carré, qui change de sens en $0$ : ici, pas
besoin de regarder le signe des nombres pour comparer leurs cubes.}
\end{exo}

\begin{exo}[Calculs et équations]
\cas{1)}
\begin{calculs}
3^{3} &= 3\times 3\times 3 = \res{27} &&\\
(-2)^{3} &= (-2)\times(-2)\times(-2) = \res{-8} &&\justif{trois facteurs négatifs : négatif}\\
\left(\frac{1}{3}\right)^{3} &= \frac{1\times 1\times 1}{3\times 3\times 3} = \res{\frac{1}{27}} &&\\
(-0{,}2)^{3} &= -0{,}2\times 0{,}2\times 0{,}2 = \res{-0{,}008} &&\justif{$2\times 2\times 2=8$ ; trois chiffres après la virgule}
\end{calculs}
\cas{2)}
\explique{L'équation $x^{3}=k$ a une seule solution, la racine cubique
$\sqrt[3]{k}$.}{méthode 9 : on isole d'abord $x^{3}$}
\cas{a)}
\explique{$x^{3}=-125$ : $x=\sqrt[3]{-125}$, soit $x=-5$. \res{$S=\ens{-5}$}.}{contrôle :
$(-5)^{3}=-125$}
\cas{b)}
\explique{$x^{3}=1$ : $x=\sqrt[3]{1}$, soit $x=1$. \res{$S=\ens{1}$}.}{une seule
solution : $-1$ ne convient pas, $(-1)^{3}=-1$}
\cas{c)}
\begin{calculs}
3x^{3}+24 &= 0 &&\\
3x^{3} &= -24 &&\justif{on retire $24$}\\
x^{3} &= -8 &&\justif{on divise par $3$}
\end{calculs}
\explique{$x=\sqrt[3]{-8}$, soit $x=-2$. \res{$S=\ens{-2}$}.}{contrôle :
$3\times(-8)+24=0$}
\cas{d)}
\begin{calculs}
x^{3} &= 4x &&\\
x^{3}-4x &= 0 &&\justif{on se ramène à « $\ldots=0$ »}\\
x\left(x^{2}-4\right) &= 0 &&\justif{facteur commun $x$}\\
x(x-2)(x+2) &= 0 &&\justif{$x^{2}-2^{2}=(x-2)(x+2)$}
\end{calculs}
\explique{Un produit est nul si l'un de ses facteurs est nul : $x=0$ ou $x=2$ ou $x=-2$.
\res{$S=\ens{-2\,;0\,;2}$}.}{trois solutions : ce n'est pas une équation $x^{3}=k$}
\rem{En d), « diviser par $x$ » ferait perdre la solution $0$ : on factorise et l'on
utilise le produit nul.}
\end{exo}

\begin{exo}[Encadrer des cubes]
\cas{1)}
\explique{$-3\le x\le 2$ donne $(-3)^{3}\le x^{3}\le 2^{3}$, soit
\res{$-27\le x^{3}\le 8$}.}{la fonction cube conserve l'ordre, même si l'intervalle
contient $0$}
\rem{Différence avec le carré : pas de piège en $0$, on prend simplement le cube des
bornes.}
\cas{2)}
\explique{$-2{,}1<-1{,}9$, donc \res{$(-2{,}1)^{3}<(-1{,}9)^{3}$}.}{la fonction cube est
croissante sur $\R$ : elle conserve l'ordre, même pour des négatifs}
\cas{3)}
\explique{Le volume d'un cube d'arête $x$ est $V=x^{3}$.}{longueur $\times$ largeur
$\times$ hauteur}
\explique{$2\le x\le 3$ donne $2^{3}\le V\le 3^{3}$, soit \res{$8\le V\le 27$} (en
cm$^{3}$).}{la fonction cube conserve l'ordre}
\explique{$x^{3}=64$ donne $x=\sqrt[3]{64}$, soit $x=4$ : une arête de
\res{$4$~cm}.}{contrôle : $4\times 4\times 4=64$}
\end{exo}

\partie{Fonction inverse}

\begin{exo}[Variations de la fonction inverse]
\cas{1)}
\begin{calculs}
\frac{1}{a}-\frac{1}{b} &= \frac{1\fois{b}}{a\fois{b}}-\frac{1\fois{a}}{b\fois{a}} &&\justif{même dénominateur $ab$}\\
\frac{1}{a}-\frac{1}{b} &= \frac{b-a}{ab} &&
\end{calculs}
\explique{$b-a>0$ (car $a<b$) et $ab>0$ (deux nombres positifs) :
$\frac{1}{a}-\frac{1}{b}>0$, donc $\frac{1}{a}>\frac{1}{b}$.}{un positif divisé par un
positif}
\explique{La fonction inverse est \res{strictement décroissante} sur
$\intoo{0}{+\infty}$.}{$a<b$ entraîne $\frac{1}{a}>\frac{1}{b}$ : elle inverse l'ordre}
\cas{2)}
\explique{Si $a<b<0$ : le calcul est le même, $b-a>0$, et $ab>0$ comme produit de deux
négatifs. Donc $\frac{1}{a}>\frac{1}{b}$.}{seul le signe de $ab$ était à revoir}
\explique{La fonction inverse est \res{strictement décroissante} sur
$\intoo{-\infty}{0}$.}{}
\cas{3)}
\begin{center}
\begin{tikzpicture}[scale=0.8]
  \tkzTabInit[lgt=1.6,espcl=2.4]{$x$/1,$\frac{1}{x}$/2}{$-\infty$,$0$,$+\infty$}
  \tkzTabVar{+/,-D+//,-/}
\end{tikzpicture}
\end{center}
\explique{La double barre sous $0$ : $0$ n'a pas d'image.}{deux flèches qui descendent,
une sur chaque intervalle}
\cas{4)}
\explique{$\dfrac{1}{-2}=-\dfrac{1}{2}$ : on a $-2<2$ et $-\frac{1}{2}<\frac{1}{2}$,
l'ordre est \emph{conservé}.}{un négatif est plus petit qu'un positif}
\explique{La fonction inverse \res{n'est pas décroissante sur $\Rs$} : elle l'est sur
chacun des deux intervalles, séparément.}{un seul contre-exemple suffit à le montrer}
\rem{On écrit « décroissante sur $\intoo{-\infty}{0}$ et sur $\intoo{0}{+\infty}$ »,
jamais « décroissante sur $\Rs$ ».}
\end{exo}

\begin{exo}[Équations et encadrements]
\cas{1) a)}
\explique{Pour $x\ne 0$ : $\frac{1}{x}=8$ donne $x=\frac{1}{8}$.
\res{$S=\ens{\frac{1}{8}}$}.}{$\frac{1}{x}=k$ donne $x=\frac{1}{k}$ ; contrôle :
$1\div\frac{1}{8}=8$}
\cas{1) b)}
\explique{$x$ est l'inverse de $-\frac{2}{3}$ : $x=-\frac{3}{2}$.
\res{$S=\ens{-\frac{3}{2}}$}.}{l'inverse d'une fraction : on échange numérateur et
dénominateur}
\cas{1) c)}
\begin{calculs}
\frac{4}{x} &= 6 &&\justif{$x\ne 0$}\\
4 &= 6x &&\justif{on multiplie par $x$}\\
x &= \frac{4}{6} &&\justif{on divise par $6$}\\
x &= \frac{2}{3} &&\justif{on simplifie par $2$}
\end{calculs}
\explique{\res{$S=\ens{\frac{2}{3}}$}.}{contrôle : $4\div\frac{2}{3}=4\times\frac{3}{2}$,
soit $6$}
\cas{1) d)}
\begin{calculs}
\frac{5}{x} &= \frac{x}{5} &&\justif{$x\ne 0$}\\
5\times 5 &= x\times x &&\justif{produit en croix}\\
x^{2} &= 25 &&
\end{calculs}
\explique{$25>0$ : \res{$S=\ens{-5\,;5}$}.}{deux solutions, toutes deux non nulles}
\cas{2)}
\explique{$2\le x\le 8$ donne $\frac{1}{8}\le\frac{1}{x}\le\frac{1}{2}$ :
\res{$\frac{1}{8}\le\frac{1}{x}\le\frac{1}{2}$}.}{nombres positifs : la fonction inverse
inverse l'ordre}
\explique{$-10\le x\le -5$ donne $\frac{1}{-5}\le\frac{1}{x}\le\frac{1}{-10}$, soit
\res{$-\frac{1}{5}\le\frac{1}{x}\le-\frac{1}{10}$}.}{nombres négatifs : l'ordre
s'inverse aussi ; contrôle : $-0{,}2\le -0{,}1$}
\cas{3)}
\explique{$0<0{,}3<0{,}4$, donc \res{$\frac{1}{0{,}3}>\frac{1}{0{,}4}$}.}{nombres
positifs : la fonction inverse inverse l'ordre}
\end{exo}

\begin{exo}[Équations et inéquations quotient]
\cas{1)}
\explique{Valeur interdite : $x-2=0$ pour $x=\res{2}$. On résout pour $x\ne 2$.}{on ne
peut pas diviser par $0$}
\begin{calculs}
\frac{x+3}{x-2} &= \frac{3}{2} &&\\
2(x+3) &= 3(x-2) &&\justif{produit en croix}\\
2x+6 &= 3x-6 &&\\
12 &= x &&\justif{on retire $2x$, on ajoute $6$}
\end{calculs}
\explique{$12\ne 2$ : la solution convient, \res{$S=\ens{12}$}.}{contrôle :
$\frac{15}{10}=\frac{3}{2}$}
\cas{2) a)}
\explique{Valeur interdite : $0$. On ne multiplie pas par $x$, dont le signe est inconnu :
on se ramène à un quotient comparé à $0$.}{comme au chapitre 1 : un seul quotient, puis un
tableau de signes}
\begin{calculs}
\frac{1}{x}-3 &= \frac{1}{x}-\frac{3\fois{x}}{1\fois{x}} &&\justif{même dénominateur $x$}\\
\frac{1}{x}-3 &= \frac{1-3x}{x} &&
\end{calculs}
\explique{On résout $\frac{1-3x}{x}\le 0$. $1-3x=0$ pour $x=\frac{1}{3}$ ; $x=0$ est la
valeur interdite.}{$1-3x$ : coefficient $-3<0$, positif puis négatif}
\begin{center}
\begin{tikzpicture}[scale=0.8]
  \tkzTabInit[lgt=2.2,espcl=2]{$x$/1,$1-3x$/1,$x$/1,$\frac{1-3x}{x}$/1.5}%
    {$-\infty$,$0$,$\frac{1}{3}$,$+\infty$}
  \tkzTabLine{,+,t,+,z,-,}
  \tkzTabLine{,-,z,+,t,+,}
  \tkzTabLine{,-,d,+,z,-,}
\end{tikzpicture}
\end{center}
\explique{\res{$S=\intoo{-\infty}{0}\cup\intfo{\frac{1}{3}}{+\infty}$}.}{signe $-$ ou
zéro ; $0$ exclu (double barre), $\frac{1}{3}$ compris (inégalité large)}
\cas{2) b)}
\explique{Valeur interdite : $1$. On se ramène à $\frac{4}{x-1}-2\ge 0$.}{on retire $2$
des deux côtés}
\begin{calculs}
\frac{4}{x-1}-2 &= \frac{4}{x-1}-\frac{2\fois{(x-1)}}{1\fois{(x-1)}} &&\justif{même dénominateur $x-1$}\\
\frac{4}{x-1}-2 &= \frac{4-2x+2}{x-1} &&\justif{$-2(x-1)=-2x+2$}\\
\frac{4}{x-1}-2 &= \frac{6-2x}{x-1} &&
\end{calculs}
\explique{$6-2x=0$ pour $x=3$ ; $x-1=0$ pour $x=1$.}{$6-2x$ : positif puis négatif}
\begin{center}
\begin{tikzpicture}[scale=0.8]
  \tkzTabInit[lgt=2.2,espcl=2]{$x$/1,$6-2x$/1,$x-1$/1,$\frac{6-2x}{x-1}$/1.5}%
    {$-\infty$,$1$,$3$,$+\infty$}
  \tkzTabLine{,+,t,+,z,-,}
  \tkzTabLine{,-,z,+,t,+,}
  \tkzTabLine{,-,d,+,z,-,}
\end{tikzpicture}
\end{center}
\explique{\res{$S=\intof{1}{3}$}.}{signe $+$ ou zéro ; $1$ exclu (valeur interdite), $3$
compris ; contrôle : $\frac{4}{3-1}=2$}
\rem{On ne multiplie jamais une inéquation par $x$ ou par $x-1$, dont le signe est
inconnu : le sens de l'inégalité pourrait changer. La valeur interdite est toujours
exclue.}
\end{exo}

\partie{Fonction racine carrée}

\begin{exo}[Se ramener au second degré]
\cas{1)}
\explique{Une solution vérifie $x-20\ge 0$, soit $x\ge 20$.}{une racine carrée est
positive, donc le membre de droite aussi}
\begin{calculs}
\sqrt{x} &= x-20 &&\justif{on recopie l'équation}\\
x &= (x-20)^{2} &&\justif{on élève au carré}\\
x &= x^{2}-40x+400 &&\justif{identité remarquable}\\
0 &= x^{2}-41x+400 &&\justif{une équation du second degré}\\
0 &= (x-16)(x-25) &&\justif{factorisation donnée}
\end{calculs}
\explique{$x=16$ ou $x=25$.}{un produit est nul si l'un de ses facteurs est nul}
\explique{$16<20$ ne convient pas ; $25$ convient. \res{$S=\ens{25}$}.}{on garde les
solutions qui vérifient $x\ge 20$ ; contrôle : $\sqrt{25}=5$ et $25-20=5$}
\cas{2)}
\explique{Une solution vérifie $x-30\ge 0$, soit $x\ge 30$.}{même démarche}
\begin{calculs}
\sqrt{x} &= x-30 &&\justif{on recopie l'équation}\\
x &= (x-30)^{2} &&\justif{on élève au carré}\\
x &= x^{2}-60x+900 &&\justif{identité remarquable}\\
0 &= x^{2}-61x+900 &&\\
0 &= (x-25)(x-36) &&\justif{factorisation donnée}
\end{calculs}
\explique{$x=25$ ou $x=36$ ; $25<30$ ne convient pas. \res{$S=\ens{36}$}.}{contrôle :
$\sqrt{36}=6$ et $36-30=6$ ; pour $25$ : $\sqrt{25}=5$ mais $25-30=-5$}
\cas{3)}
\begin{minipage}[t]{0.56\linewidth}\raggedright
\vspace{0pt}
À la calculatrice, on trace la courbe de $x\mapsto\sqrt{x+3}$ et la droite d'équation
$y=x+1$. Les solutions de $(E)$ sont les abscisses de leurs points communs.\par
Il n'y en a qu'un, apparemment $(1\,;2)$ : on conjecture que $(E)$ a \res{une seule
solution, $1$}.\par
\gris{La courbe part du point $(-3\,;0)$ : $\sqrt{x+3}$ n'existe que pour $x\ge -3$.}
\end{minipage}\hfill
\begin{minipage}[t]{0.40\linewidth}\centering
\vspace{0pt}
\begin{tikzpicture}[x=0.45cm,y=0.45cm]
  \repere{-4}{-3}{4}{5}
  \gradx{-3,-2,-1,1,2,3} \grady{-2,-1,1,2,3,4}
  \draw[coulCourbe,line width=1.1pt,domain=0:2.6458,samples=60] plot ({\x*\x-3},\x);
  \draw[coulCourbeC,line width=1.1pt] (-4,-3) -- (4,5);
  \draw[lecture] (1,2) -- (1,0);
  \path[coulProp] (1,2) node[croix]{};
  \node[coulCourbe,font=\footnotesize] at (3,1.6) {$y=\sqrt{x+3}$};
  \node[coulCourbeC,font=\footnotesize] at (2.3,4.6) {$y=x+1$};
\end{tikzpicture}
\end{minipage}
\cas{4)}
\explique{Une racine carrée est positive : si $x$ est solution, $x+1=\sqrt{x+3}\ge 0$,
donc \res{$x+1\ge 0$}.}{c'est la condition, à écrire avant d'élever au carré}
\begin{calculs}
\sqrt{x+3} &= x+1 &&\justif{on recopie l'équation}\\
x+3 &= (x+1)^{2} &&\justif{on élève au carré}\\
x+3 &= x^{2}+2x+1 &&\justif{identité remarquable}\\
0 &= x^{2}+x-2 &&\justif{on retire $x+3$ des deux côtés}
\end{calculs}
\explique{Toute solution de $(E)$ vérifie donc \res{$x^{2}+x-2=0$}.}{}
\cas{5)}
\begin{calculs}
(x+2)(x-1) &= x^{2}-x+2x-2 &&\justif{double distributivité}\\
(x+2)(x-1) &= x^{2}+x-2 &&
\end{calculs}
\explique{L'équation $x^{2}+x-2=0$ s'écrit $(x+2)(x-1)=0$ : $x=-2$ ou $x=1$.}{produit
nul}
\explique{$-2+1=-1<0$ : $-2$ ne convient pas. Pour $1$ : $\sqrt{1+3}=2$ et $1+1=2$.
\res{$S=\ens{1}$}.}{on revient à la condition $x+1\ge 0$ ; la conjecture du 3) est
démontrée}
\cas{6)}
\explique{Une solution vérifie $x\ge 0$ (pour $\sqrt{x}$) et $x-3\ge 0$, soit
$x\ge 3$.}{la seconde condition contient la première}
\begin{calculs}
2\sqrt{x} &= x-3 &&\justif{on recopie l'équation}\\
4x &= (x-3)^{2} &&\justif{on élève au carré : $\left(2\sqrt{x}\right)^{2}=4x$}\\
4x &= x^{2}-6x+9 &&\justif{identité remarquable}\\
0 &= x^{2}-10x+9 &&\justif{on retire $4x$}
\end{calculs}
\explique{On vérifie la forme proposée, puis on factorise avec
$A^{2}-B^{2}=(A-B)(A+B)$.}{ici $A=x-5$ et $B=4$}
\begin{calculs}
(x-5)^{2}-4^{2} &= x^{2}-10x+25-16 &&\\
(x-5)^{2}-4^{2} &= x^{2}-10x+9 &&\justif{la forme est la bonne}
\end{calculs}
\begin{calculs}
x^{2}-10x+9 &= (x-5-4)(x-5+4) &&\\
x^{2}-10x+9 &= (x-9)(x-1) &&
\end{calculs}
\explique{$x=9$ ou $x=1$ ; $1<3$ ne convient pas. Pour $9$ : $2\sqrt{9}=6$ et $9-3=6$.
\res{$S=\ens{9}$}.}{pour $1$ : $2\sqrt{1}=2$ mais $1-3=-2$, une fausse solution}
\rem{Élever au carré fait apparaître de fausses solutions : on revient toujours à la
condition ($x\ge 20$ ; $x+1\ge 0$ ; $x\ge 3$) et on contrôle dans l'équation de départ.}
\end{exo}

\begin{exo}[Équations avec des racines carrées]
\explique{Pour chaque équation : la condition d'existence, puis on isole la racine et on
élève au carré, puis on contrôle.}{méthode 11}
\cas{a)}
\begin{calculs}
\sqrt{x} &= 9 &&\justif{il faut $x\ge 0$}\\
x &= 81 &&\justif{on élève au carré}
\end{calculs}
\explique{\res{$S=\ens{81}$}.}{contrôle : $\sqrt{81}=9$}
\cas{b)}
\explique{Une racine carrée est positive, jamais égale à $-3$ : \res{$S=\varnothing$}.}{pas
de calcul : élever au carré donnerait $x=9$, une fausse solution}
\cas{c)}
\begin{calculs}
\sqrt{x-4} &= 5 &&\justif{il faut $x-4\ge 0$, soit $x\ge 4$}\\
x-4 &= 25 &&\justif{on élève au carré}\\
x &= 29 &&\justif{$29\ge 4$}
\end{calculs}
\explique{\res{$S=\ens{29}$}.}{contrôle : $\sqrt{25}=5$}
\cas{d)}
\begin{calculs}
3\sqrt{2x+1} &= 15 &&\justif{il faut $2x+1\ge 0$, soit $x\ge -\frac{1}{2}$}\\
\sqrt{2x+1} &= 5 &&\justif{on isole la racine : on divise par $3$}\\
2x+1 &= 25 &&\justif{on élève au carré}\\
2x &= 24 &&\\
x &= 12 &&\justif{$12\ge -\frac{1}{2}$}
\end{calculs}
\explique{\res{$S=\ens{12}$}.}{contrôle : $3\sqrt{25}=3\times 5$, soit $15$}
\end{exo}

\begin{exo}[Calculer avec des racines carrées]
\explique{Règles utilisées : $\sqrt{a\times b}=\sqrt{a}\times\sqrt{b}$ et
$\left(\sqrt{a}\right)^{2}=a$ pour $a\ge 0$ et $b\ge 0$.}{on fait apparaître un carré
parfait sous la racine}
\cas{1) a)}
\begin{calculs}
\sqrt{18}-\sqrt{2} &= \sqrt{9\times 2}-\sqrt{2} &&\justif{$9$ est un carré parfait}\\
\sqrt{18}-\sqrt{2} &= 3\sqrt{2}-\sqrt{2} &&\justif{$\sqrt{9}\times\sqrt{2}$}\\
\sqrt{18}-\sqrt{2} &= \res{2\sqrt{2}} &&\justif{$3$ fois $\sqrt{2}$, moins $1$ fois $\sqrt{2}$}
\end{calculs}
\cas{1) b)}
\begin{calculs}
\left(\sqrt{5}-\sqrt{2}\right)\left(\sqrt{5}+\sqrt{2}\right) &= \left(\sqrt{5}\right)^{2}-\left(\sqrt{2}\right)^{2} &&\justif{$(a-b)(a+b)=a^{2}-b^{2}$}\\
\left(\sqrt{5}-\sqrt{2}\right)\left(\sqrt{5}+\sqrt{2}\right) &= 5-2 &&\\
\left(\sqrt{5}-\sqrt{2}\right)\left(\sqrt{5}+\sqrt{2}\right) &= \res{3} &&
\end{calculs}
\cas{1) c)}
\begin{calculs}
\left(\sqrt{3}+\sqrt{6}\right)^{2} &= \left(\sqrt{3}\right)^{2}+2\sqrt{3}\sqrt{6}+\left(\sqrt{6}\right)^{2} &&\justif{$(a+b)^{2}=a^{2}+2ab+b^{2}$}\\
\left(\sqrt{3}+\sqrt{6}\right)^{2} &= 3+2\sqrt{18}+6 &&\justif{$\sqrt{3}\times\sqrt{6}=\sqrt{18}$}\\
\left(\sqrt{3}+\sqrt{6}\right)^{2} &= 9+2\times 3\sqrt{2} &&\justif{$\sqrt{18}=3\sqrt{2}$, vu en a)}\\
\left(\sqrt{3}+\sqrt{6}\right)^{2} &= \res{9+6\sqrt{2}} &&
\end{calculs}
\rem{Erreur fréquente : $\left(\sqrt{3}+\sqrt{6}\right)^{2}=3+6$. On oublie le double
produit.}
\cas{2)}
\begin{calculs}
A^{2} &= \left(3-\sqrt{10}\right)^{2} &&\\
A^{2} &= 3^{2}-2\times 3\sqrt{10}+\left(\sqrt{10}\right)^{2} &&\justif{$(a-b)^{2}=a^{2}-2ab+b^{2}$}\\
A^{2} &= 9-6\sqrt{10}+10 &&\\
A^{2} &= \res{19-6\sqrt{10}} &&
\end{calculs}
\begin{calculs}
B^{2} &= \left(\sqrt{19-6\sqrt{10}}\right)^{2} &&\\
B^{2} &= \res{19-6\sqrt{10}} &&\justif{$\left(\sqrt{a}\right)^{2}=a$}
\end{calculs}
\explique{$A^{2}=B^{2}$ donne seulement $A=B$ ou $A=-B$ : il faut regarder les
signes.}{deux nombres opposés ont le même carré}
\explique{$\sqrt{10}>\sqrt{9}$, soit $\sqrt{10}>3$ : donc $A=3-\sqrt{10}<0$. Or $B\ge 0$,
comme toute racine carrée.}{la fonction racine carrée est croissante}
\explique{Donc \res{$A\ne B$} : en fait, $B=-A=\sqrt{10}-3$.}{$A$ est négatif, $B$
positif}
\cas{3) a)}
\begin{calculs}
\frac{1}{\sqrt{3}} &= \frac{1\fois{\sqrt{3}}}{\sqrt{3}\fois{\sqrt{3}}} &&\justif{on multiplie en haut et en bas par $\sqrt{3}$}\\
\frac{1}{\sqrt{3}} &= \res{\frac{\sqrt{3}}{3}} &&\justif{$\sqrt{3}\times\sqrt{3}=3$ : plus de racine au dénominateur}
\end{calculs}
\cas{3) b)}
\begin{calculs}
\frac{1}{\sqrt{7}-2} &= \frac{1\fois{\left(\sqrt{7}+2\right)}}{\left(\sqrt{7}-2\right)\fois{\left(\sqrt{7}+2\right)}} &&\justif{on multiplie en haut et en bas par $\sqrt{7}+2$}\\
\frac{1}{\sqrt{7}-2} &= \frac{\sqrt{7}+2}{\left(\sqrt{7}\right)^{2}-2^{2}} &&\justif{$(a-b)(a+b)=a^{2}-b^{2}$}\\
\frac{1}{\sqrt{7}-2} &= \frac{\sqrt{7}+2}{7-4} &&\\
\frac{1}{\sqrt{7}-2} &= \res{\frac{\sqrt{7}+2}{3}} &&
\end{calculs}
\explique{L'égalité est justifiée.}{le facteur $\sqrt{7}+2$ est choisi pour faire
apparaître $a^{2}-b^{2}$, sans racine}
\cas{3) c)}
\begin{calculs}
\sqrt{50}-\sqrt{32}+\sqrt{2} &= \sqrt{25\times 2}-\sqrt{16\times 2}+\sqrt{2} &&\justif{carrés parfaits $25$ et $16$}\\
\sqrt{50}-\sqrt{32}+\sqrt{2} &= 5\sqrt{2}-4\sqrt{2}+\sqrt{2} &&\\
\sqrt{50}-\sqrt{32}+\sqrt{2} &= \res{2\sqrt{2}} &&\justif{$5-4+1=2$}
\end{calculs}
\cas{4)}
\begin{calculs}
\left(\sqrt{6}+\sqrt{7}\right)\left(\sqrt{6}-\sqrt{7}\right) &= \left(\sqrt{6}\right)^{2}-\left(\sqrt{7}\right)^{2} &&\justif{$(a+b)(a-b)=a^{2}-b^{2}$}\\
\left(\sqrt{6}+\sqrt{7}\right)\left(\sqrt{6}-\sqrt{7}\right) &= 6-7 &&\\
\left(\sqrt{6}+\sqrt{7}\right)\left(\sqrt{6}-\sqrt{7}\right) &= \res{-1} &&
\end{calculs}
\explique{Le produit est négatif, et $\sqrt{6}+\sqrt{7}>0$ : donc $\sqrt{6}-\sqrt{7}<0$,
c'est-à-dire \res{$\sqrt{6}<\sqrt{7}$}.}{un produit négatif a deux facteurs de signes
contraires}
\end{exo}

\partie{Fonction valeur absolue}

\begin{exo}[Étude de $x\mapsto|x+3|$]
\cas{1)}
\begin{center}
{\renewcommand{\arraystretch}{1.3}\setlength{\tabcolsep}{5pt}
\begin{tabular}{|c|*{7}{c|}}
\hline
$x$ & $-6$ & $-5$ & $-4$ & $-3$ & $-2$ & $-1$ & $0$\\
\hline
$f(x)$ & $3$ & $2$ & $1$ & $0$ & $1$ & $2$ & $3$\\
\hline
\end{tabular}}
\end{center}
\explique{Par exemple, $f(-6)=|-6+3|=|-3|$, soit $3$.}{on calcule d'abord dans la valeur
absolue, puis on prend la distance à $0$}
\begin{minipage}[t]{0.58\linewidth}\raggedright
\vspace{0pt}
\cas{2)}
\explique{On place les points du tableau, puis on les relie : deux demi-droites en
« V ».}{figure ci-contre}
\cas{3)}
\begin{center}
\begin{tikzpicture}[scale=0.75]
  \tkzTabInit[lgt=1.8,espcl=2.2]{$x$/1,$|x+3|$/2}{$-\infty$,$-3$,$+\infty$}
  \tkzTabVar{+/,-/$0$,+/}
\end{tikzpicture}
\end{center}
\explique{Le minimum de $f$ est \res{$0$, atteint en $-3$}.}{décroissante avant le
sommet, croissante après}
\cas{4)}
\explique{C'est le « V » de la valeur absolue, \res{décalé de $3$ unités vers la
gauche}, de sommet $(-3\,;0)$.}{le sommet est là où $x+3=0$}
\end{minipage}\hfill
\begin{minipage}[t]{0.38\linewidth}\centering
\vspace{0pt}
\begin{tikzpicture}[x=0.5cm,y=0.5cm]
  \repere{-7}{-1}{1}{4}
  \gradx{-6,-5,-4,-3,-2,-1} \grady{1,2,3}
  \draw[coulCourbe,line width=1.1pt] (-6.5,3.5) -- (-3,0) -- (0.5,3.5);
  \draw[coulCourbeC,thick] (-7,2) -- (1,2);
  \draw[lecture] (-5,0) -- (-5,2) (-1,0) -- (-1,2);
  \path[coulCourbe] (-6,3) node[croix]{} (-5,2) node[croix]{} (-4,1) node[croix]{}
    (-3,0) node[croix]{} (-2,1) node[croix]{} (-1,2) node[croix]{} (0,3) node[croix]{};
  \node[coulCourbeC,font=\footnotesize] at (-6.2,2.5) {$y=2$};
  \node[coulCourbe,font=\footnotesize] at (-0.9,3.6) {$\Cf$};
\end{tikzpicture}
\end{minipage}
\cas{5)}
\begin{calculs}
|x+3| &= 2 &&\justif{on recopie l'équation}\\
x+3 &= 2 \quad\text{ou}\quad x+3=-2 &&\justif{$|A|=2$ : $A=2$ ou $A=-2$}\\
x &= -1 \quad\text{ou}\quad x=-5 &&
\end{calculs}
\explique{\res{$S=\ens{-5\,;-1}$}.}{les deux nombres à distance $2$ de $-3$}
\explique{Une valeur absolue est positive : $|x+3|=-5$ n'a pas de solution,
\res{$S=\varnothing$}.}{inutile de calculer}
\cas{6)}
\explique{La droite d'équation $y=2$ coupe $\Cf$ en deux points, d'abscisses $-5$ et
$-1$.}{on retrouve $S=\ens{-5\,;-1}$}
\explique{$\Cf$ n'a aucun point sous l'axe des abscisses, donc aucun point d'ordonnée
$-5$.}{on retrouve $S=\varnothing$}
\end{exo}

\begin{exo}[Calculs et distances]
\explique{$|x|=x$ si $x\ge 0$ et $|x|=-x$ si $x<0$ : c'est la distance de $x$ à
$0$.}{définition 9 ; on calcule d'abord ce qui est entre les barres}
\cas{1)}
\begin{calculs}
|-12| &= \res{12} &&\justif{$-12<0$ : on prend l'opposé}\\
|3-10| &= |-7| &&\justif{on calcule d'abord dans les barres}\\
|3-10| &= \res{7} &&\\
|-4|\times|-5| &= 4\times 5 &&\\
|-4|\times|-5| &= \res{20} &&\\
\left|\sqrt{2}-2\right| &= -\left(\sqrt{2}-2\right) &&\justif{$\sqrt{2}\approx 1{,}41<2$ : $\sqrt{2}-2<0$, on prend l'opposé}\\
\left|\sqrt{2}-2\right| &= \res{2-\sqrt{2}} &&\\
|-6|-|2| &= 6-2 &&\\
|-6|-|2| &= \res{4} &&
\end{calculs}
\rem{En d), on cherche d'abord le signe de $\sqrt{2}-2$ : $|A|=-A$ lorsque $A<0$.
Erreur fréquente : écrire $\sqrt{2}-2$, qui est négatif.}
\cas{2)}
\begin{calculs}
\sqrt{(-8)^{2}} &= \sqrt{64} &&\justif{le carré d'abord}\\
\sqrt{(-8)^{2}} &= \res{8} &&\\
\sqrt{(-0{,}5)^{2}} &= \sqrt{0{,}25} &&\\
\sqrt{(-0{,}5)^{2}} &= \res{0{,}5} &&
\end{calculs}
\explique{On ne retrouve pas $-8$ mais son opposé : en général,
\res{$\sqrt{a^{2}}=|a|$}.}{$\sqrt{a^{2}}=a$ est faux pour $a<0$ : une racine carrée est
positive}
\cas{3)}
\begin{calculs}
|3-(-6)| &= |9| &&\justif{distance entre $a$ et $b$ : $|b-a|$}\\
|3-(-6)| &= \res{9} &&\\
|-1-7| &= |-8| &&\\
|-1-7| &= \res{8} &&
\end{calculs}
\explique{La distance entre $-6$ et $3$ est $9$ ; entre $7$ et $-1$, elle est $8$.}{l'ordre
des nombres ne compte pas : $|7-(-1)|=8$ aussi}
\end{exo}

\begin{exo}[Équations avec des valeurs absolues]
\explique{Si $r>0$, $|A|=r$ équivaut à $A=r$ ou $A=-r$ ; si $r<0$, pas de
solution.}{propriété 4}
\cas{a)}
\explique{$|x|=9$ : $x=9$ ou $x=-9$, \res{$S=\ens{-9\,;9}$}.}{les deux nombres à
distance $9$ de $0$}
\cas{b)}
\explique{Une valeur absolue est positive : $|x|=-3$ n'a pas de solution,
\res{$S=\varnothing$}.}{$r=-3<0$}
\cas{c)}
\begin{calculs}
|x+2| &= 6 &&\justif{on recopie l'équation}\\
x+2 &= 6 \quad\text{ou}\quad x+2=-6 &&\justif{$A=6$ ou $A=-6$}\\
x &= 4 \quad\text{ou}\quad x=-8 &&\justif{on retire $2$}
\end{calculs}
\explique{\res{$S=\ens{-8\,;4}$}.}{les deux nombres à distance $6$ de $-2$}
\cas{d)}
\begin{calculs}
|3x-1| &= 8 &&\\
3x-1 &= 8 \quad\text{ou}\quad 3x-1=-8 &&\\
3x &= 9 \quad\text{ou}\quad 3x=-7 &&\justif{on ajoute $1$}\\
x &= 3 \quad\text{ou}\quad x=-\frac{7}{3} &&\justif{on divise par $3$}
\end{calculs}
\explique{\res{$S=\ens{-\frac{7}{3}\,;3}$}.}{contrôle : $|9-1|=8$ et $|-7-1|=8$}
\end{exo}

\begin{exo}[Valeur absolue égale à une expression]
\cas{1)}
\explique{Une valeur absolue est positive : si $x$ est solution, $x+1=|2x-4|\ge 0$, donc
$x+1\ge 0$, soit \res{$x\ge -1$}.}{même idée que pour $\sqrt{x}=x-k$ (ex. 27) : le membre
de droite ne peut pas être négatif}
\cas{2)}
\begin{calculs}
2x-4 &= x+1 &&\justif{premier cas : $A=x+1$}\\
x &= 5 &&\justif{on retire $x$, on ajoute $4$}
\end{calculs}
\begin{calculs}
2x-4 &= -(x+1) &&\justif{second cas : $A=-(x+1)$}\\
2x-4 &= -x-1 &&\\
3x &= 3 &&\justif{on ajoute $x$ et $4$}\\
x &= 1 &&
\end{calculs}
\cas{3)}
\explique{Il faut $x\ge -1$ : $5$ et $1$ conviennent tous les deux.}{on trie les
solutions avec la condition du 1)}
\explique{Contrôle : $|10-4|=6$ et $5+1=6$ ; $|2-4|=2$ et $1+1=2$. \res{$S=\ens{1\,;5}$}.}{on
vérifie dans l'équation de départ}
\cas{4)}
\explique{$|x+3|=2x$ : il faut $2x\ge 0$, soit $x\ge 0$.}{la condition d'abord}
\begin{calculs}
x+3 &= 2x &&\justif{premier cas}\\
x &= 3 &&
\end{calculs}
\begin{calculs}
x+3 &= -2x &&\justif{second cas}\\
3x &= -3 &&\\
x &= -1 &&
\end{calculs}
\explique{$-1<0$ ne convient pas ; contrôle pour $3$ : $|6|=6$ et $2\times 3=6$.
\res{$S=\ens{3}$}.}{pour $-1$ : $|2|=2$ mais $2\times(-1)=-2$}
\explique{$|3x-1|=x-5$ : il faut $x-5\ge 0$, soit $x\ge 5$.}{}
\begin{calculs}
3x-1 &= x-5 &&\justif{premier cas}\\
2x &= -4 &&\\
x &= -2 &&
\end{calculs}
\begin{calculs}
3x-1 &= -(x-5) &&\justif{second cas}\\
3x-1 &= -x+5 &&\\
4x &= 6 &&\\
x &= \frac{3}{2} &&
\end{calculs}
\explique{Aucune des deux ne vérifie $x\ge 5$ : \res{$S=\varnothing$}.}{deux fausses
solutions}
\rem{Comme pour $\sqrt{x}=x-k$, on revient toujours à la condition : sans elle, on garde
de fausses solutions.}
\end{exo}

\partie{Fonctions affines}

\begin{exo}[Lire l'expression d'une fonction affine]
\begin{minipage}[t]{0.56\linewidth}\raggedright
\vspace{0pt}
On lit l'ordonnée à l'origine $b$ (où la droite coupe l'axe des ordonnées), puis le
coefficient directeur $a$ avec deux points à coordonnées entières.\par
\gris{L'expression est $f(x)=ax+b$ ; la droite a pour équation $y=ax+b$.}

\cas{$(d_{1})$}
Elle passe par $(0\,;1)$ : $b=1$. De $(0\,;1)$ à $(1\,;-1)$, $x$ augmente de $1$ et
$y$ baisse de $2$ : $a=-2$.\par
L'expression est \res{$-2x+1$}.\par
\gris{Contrôle : la droite passe par $(2\,;-3)$ et $-2\times 2+1=-3$.}

\cas{$(d_{2})$}
Elle est horizontale, à la hauteur $3$ : $a=0$ et $b=3$.\par
L'expression est \res{$3$} : une fonction constante.

\cas{$(d_{3})$}
Elle passe par $(0\,;-2)$ : $b=-2$. De $(0\,;-2)$ à $(3\,;0)$, $x$ augmente de $3$ et
$y$ de $2$ :
\[
a=\frac{0-(-2)}{3-0}=\frac{2}{3}
\]
L'expression est \res{$\frac{2}{3}x-2$}.\par
\gris{Quand $x$ augmente de $1$, $y$ augmente de $\frac{2}{3}$ : cela ne tombe pas sur
le quadrillage, d'où les deux points à coordonnées entières.}
\end{minipage}\hfill
\begin{minipage}[t]{0.40\linewidth}\centering
\vspace{0pt}
\begin{tikzpicture}[x=0.5cm,y=0.5cm]
  \repere{-5}{-4}{5}{4}
  \gradx{-4,-3,-2,-1,1,2,3,4} \grady{-3,-2,-1,1,2,3}
  \begin{scope}
  \clip (-5,-4) rectangle (5,4);
  \draw[coulCourbe,line width=1.1pt] (-1.5,4) -- (2.5,-4);
  \draw[coulCourbeB,line width=1.1pt] (-5,3) -- (5,3);
  \draw[coulCourbeC,line width=1.1pt] (-3,-4) -- (5,{4/3});
  \end{scope}
  \draw[lecture] (0,1) -- (1,1) -- (1,-1);
  \node[coulProp,font=\scriptsize,above] at (0.5,1) {$+1$};
  \node[coulProp,font=\scriptsize,right] at (1,0) {$-2$};
  \draw[lecture] (0,-2) -- (3,-2) -- (3,0);
  \node[coulProp,font=\scriptsize,below] at (1.5,-2) {$+3$};
  \node[coulProp,font=\scriptsize,right] at (3,-1) {$+2$};
  \path[coulProp] (0,1) node[croix]{} (1,-1) node[croix]{} (0,3) node[croix]{}
    (0,-2) node[croix]{} (3,0) node[croix]{};
  \node[coulCourbe,font=\footnotesize] at (-0.3,3.6) {$(d_{1})$};
  \node[coulCourbeB,font=\footnotesize] at (-4.1,3.5) {$(d_{2})$};
  \node[coulCourbeC,font=\footnotesize] at (4.2,0.4) {$(d_{3})$};
\end{tikzpicture}
\end{minipage}
\rem{$b$ se lit sur l'axe des ordonnées, pas sur l'axe des abscisses. Erreur fréquente
pour $(d_{3})$ : prendre $3$, l'abscisse du point où elle coupe l'axe des abscisses.}
\end{exo}

\begin{exo}[Équation réduite et alignement]
\explique{$x_{A}\ne x_{B}$ : la droite $(AB)$ a une équation réduite
$y=ax+b$.}{méthode 13}
\cas{1)}
\begin{calculs}
a &= \frac{y_{B}-y_{A}}{x_{B}-x_{A}} &&\justif{coefficient directeur}\\
a &= \frac{-2-4}{2-(-1)} &&\justif{on remplace, nombres négatifs entre parenthèses}\\
a &= \frac{-6}{3} &&\\
a &= \res{-2} &&
\end{calculs}
\cas{2)}
\begin{calculs}
4 &= -2\times(-1)+b &&\justif{$A$ est sur la droite : $y_{A}=a\,x_{A}+b$}\\
4 &= 2+b &&\\
b &= 2 &&\justif{on retire $2$}
\end{calculs}
\explique{L'équation réduite de $(AB)$ est \res{$y=-2x+2$}.}{contrôle avec $B$ :
$-2\times 2+2=-2$, soit $y_{B}$}
\cas{3)}
\explique{$C$ est aligné avec $A$ et $B$ lorsqu'il est sur la droite $(AB)$ : on
remplace $x$ par $x_{C}$ dans l'équation et on compare à $y_{C}$.}{même chose pour $D$}
\begin{calculs}
-2\times 5+2 &= -8 &&\justif{c'est $y_{C}$}\\
-2\times(-3)+2 &= 8 &&\justif{et $y_{D}=9$}
\end{calculs}
\explique{\res{$C$ est aligné} avec $A$ et $B$ ; $8\ne 9$ : \res{$D$ n'est pas
aligné} avec $A$ et $B$.}{les coordonnées de $C$ vérifient l'équation, pas celles de $D$}
\end{exo}

\begin{exo}[Intersection de deux droites]
\explique{Un point commun a la même ordonnée sur les deux droites : on résout
« $y$ de $(d)$ $=$ $y$ de $(d')$ », puis on calcule l'ordonnée.}{l'abscisse d'abord,
l'ordonnée ensuite}
\cas{1)}
\begin{calculs}
3x-5 &= -2x+5 &&\\
5x &= 10 &&\justif{on ajoute $2x$ et $5$}\\
x &= 2 &&\\
y &= 3\times 2-5 = 1 &&\justif{avec l'équation de $(d)$}
\end{calculs}
\explique{Le point d'intersection est \res{$(2\,;1)$}.}{contrôle avec $(d')$ :
$-2\times 2+5=1$}
\cas{2)}
\explique{On cherche d'abord l'équation réduite de $(d)$. $A(0\,;-1)$ est sur l'axe des
ordonnées : $b=-1$.}{l'ordonnée à l'origine se lit directement}
\begin{calculs}
a &= \frac{5-(-1)}{3-0} &&\justif{$\frac{y_{B}-y_{A}}{x_{B}-x_{A}}$}\\
a &= \frac{6}{3} = 2 &&
\end{calculs}
\explique{$(d) : y=2x-1$.}{contrôle avec $B$ : $2\times 3-1=5$}
\begin{calculs}
2x-1 &= -x+5 &&\justif{même ordonnée sur $(d)$ et $(d')$}\\
3x &= 6 &&\justif{on ajoute $x$ et $1$}\\
x &= 2 &&\\
y &= 2\times 2-1 = 3 &&
\end{calculs}
\explique{Le point d'intersection est \res{$(2\,;3)$}.}{contrôle avec $(d')$ : $-2+5=3$}
\end{exo}

\partie{Prolongement : vers la dérivation}

\begin{exo}[Calculer des taux de variation]
\explique{Taux de variation entre $a$ et $b$ : $T=\dfrac{f(b)-f(a)}{b-a}$.}{définition 11 ;
on calcule d'abord les deux images}
\cas{1)}
\explique{$f(0)=0$ et $f(3)=9+6$, soit $15$.}{$f(x)=x^{2}+2x$}
\begin{calculs}
T &= \frac{15-0}{3-0} &&\\
T &= \res{5} &&
\end{calculs}
\explique{$f(-4)=16-8$, soit $8$ ; $f(-1)=1-2$, soit $-1$.}{$(-4)^{2}=16$ ;
$2\times(-4)=-8$}
\begin{calculs}
T &= \frac{-1-8}{-1-(-4)} &&\justif{$a=-4$ et $b=-1$}\\
T &= \frac{-9}{3} &&\\
T &= \res{-3} &&
\end{calculs}
\cas{2)}
\explique{$g(-1)=2+7$, soit $9$ ; $g(4)=-8+7$, soit $-1$.}{$g(x)=-2x+7$}
\begin{calculs}
T &= \frac{-1-9}{4-(-1)} &&\\
T &= \frac{-10}{5} = \res{-2} &&
\end{calculs}
\explique{$g(2)=-4+7$, soit $3$ ; $g(10)=-20+7$, soit $-13$.}{}
\begin{calculs}
T &= \frac{-13-3}{10-2} &&\\
T &= \frac{-16}{8} = \res{-2} &&
\end{calculs}
\explique{On trouve \res{le même nombre}, $-2$ : le coefficient directeur de la droite
qui représente $g$.}{le taux ne dépend pas des deux nombres choisis}
\cas{3)}
\begin{calculs}
T &= \frac{(mb+p)-(ma+p)}{b-a} &&\justif{avec $a\ne b$}\\
T &= \frac{mb-ma}{b-a} &&\justif{les $p$ s'annulent}\\
T &= \frac{m(b-a)}{b-a} &&\justif{facteur commun $m$}\\
T &= \res{m} &&\justif{on simplifie par $b-a\ne 0$}
\end{calculs}
\explique{La courbe d'une fonction affine est une droite : toute sécante est cette droite
elle-même, de coefficient directeur $m$.}{le taux de variation est le coefficient
directeur de la sécante}
\end{exo}

\begin{exo}[Taux et sens de variation]
\cas{1)}
\begin{calculs}
T &= \frac{g(b)-g(a)}{b-a} &&\\
T &= \frac{\left(b^{2}-4b\right)-\left(a^{2}-4a\right)}{b-a} &&\justif{parenthèses autour de $g(a)$}\\
T &= \frac{b^{2}-a^{2}-4b+4a}{b-a} &&\\
T &= \frac{(b-a)(b+a)-4(b-a)}{b-a} &&\justif{$b^{2}-a^{2}=(b-a)(b+a)$}\\
T &= \frac{(b-a)(b+a-4)}{b-a} &&\justif{facteur commun $b-a$}\\
T &= \res{a+b-4} &&\justif{on simplifie par $b-a\ne 0$}
\end{calculs}
\explique{Sur $\intfo{2}{+\infty}$ : $a\ge 2$ et $b\ge 2$, distincts, donc l'un dépasse
$2$ et $a+b>4$. Alors $T>0$ : $g$ est \res{croissante} sur $\intfo{2}{+\infty}$.}{taux
toujours positif : croissante (propriété 7)}
\explique{Sur $\intof{-\infty}{2}$ : de même $a+b<4$, $T<0$ : $g$ est \res{décroissante}
sur $\intof{-\infty}{2}$.}{taux toujours négatif : décroissante}
\cas{2)}
\begin{calculs}
T &= \frac{\frac{1}{b}-\frac{1}{a}}{b-a} &&\justif{$h(b)-h(a)$ au numérateur}\\
T &= \left(\frac{1\fois{a}}{b\fois{a}}-\frac{1\fois{b}}{a\fois{b}}\right)\times\frac{1}{b-a} &&\justif{diviser par $b-a$, c'est multiplier par $\frac{1}{b-a}$}\\
T &= \frac{a-b}{ab}\times\frac{1}{b-a} &&\\
T &= \frac{-(b-a)}{ab(b-a)} &&\justif{$a-b=-(b-a)$}\\
T &= \res{-\frac{1}{ab}} &&\justif{on simplifie par $b-a\ne 0$}
\end{calculs}
\explique{Sur $\intoo{0}{+\infty}$ : $ab>0$, donc $T<0$ : $h$ est \res{décroissante} sur
$\intoo{0}{+\infty}$.}{on retrouve la partie VII}
\end{exo}

\begin{exo}[Vitesse moyenne d'une chute]
\explique{Vitesse moyenne entre $t_{1}$ et $t_{2}$ : $\dfrac{d(t_{2})-d(t_{1})}{t_{2}-t_{1}}$,
en mètres par seconde.}{distance parcourue divisée par la durée}
\cas{1)}
\explique{$d(1)=5\times 1^{2}$, soit $5$ ; $d(3)=5\times 3^{2}$, soit $45$.}{}
\begin{calculs}
v &= \frac{45-5}{3-1} &&\\
v &= \frac{40}{2} = \res{20} &&\justif{en m/s}
\end{calculs}
\cas{2)}
\explique{$d(2)=5\times 4$, soit $20$ ; $d(2{,}5)=5\times 6{,}25$, soit $31{,}25$.}{}
\begin{calculs}
v &= \frac{31{,}25-20}{2{,}5-2} &&\\
v &= \frac{11{,}25}{0{,}5} = \res{22{,}5} &&\justif{en m/s}
\end{calculs}
\explique{$d(2{,}1)=5\times 4{,}41$, soit $22{,}05$.}{$2{,}1^{2}=4{,}41$}
\begin{calculs}
v &= \frac{22{,}05-20}{2{,}1-2} &&\\
v &= \frac{2{,}05}{0{,}1} = \res{20{,}5} &&\justif{en m/s}
\end{calculs}
\cas{3)}
\explique{Quand le second instant se rapproche de $2$, la vitesse moyenne se rapproche de
\res{$20$~m/s} : c'est la vitesse de la bille à l'instant $t=2$.}{$22{,}5$ puis
$20{,}5$ : on conjecture $20$ ; ce sera le \emph{nombre dérivé} de $d$ en $2$ (chapitre
4)}
\end{exo}

\vfill

\end{document}
